\chapter{Itapúa}\label{dept:Itapúa}
\mapaparaguai{7}{Itapúa}
\vfill\clearpage
\section{Itapua}

\begin{itemize}
\tightlist
\item
  Encarnacion e a capital, concentrando o melhor polo de saude e
  servicos do sul do pais.
\item
  Solo basaltico vermelho de alta fertilidade confirma vocacao
  agroindustrial de exportacao.
\item
  Lacuna de solar/pluviometria da capital fechada com dados NASA POWER
  (2001-2020).
\end{itemize}

\subsection{Ficha climática de Encarnacion (capital
departamental)}

Fonte: NASA POWER Climatology API, período 2001-2020, coordenadas
-27.33°S / -55.87°W.

\textbf{Irradiância solar global --- (kWh/m²/dia):}

{\def\LTcaptype{none} % do not increment counter
\begin{longtable}[]{@{}
  >{\raggedright\arraybackslash}p{(\linewidth - 22\tabcolsep) * \real{0.0833}}
  >{\raggedright\arraybackslash}p{(\linewidth - 22\tabcolsep) * \real{0.0833}}
  >{\raggedright\arraybackslash}p{(\linewidth - 22\tabcolsep) * \real{0.0833}}
  >{\raggedright\arraybackslash}p{(\linewidth - 22\tabcolsep) * \real{0.0833}}
  >{\raggedright\arraybackslash}p{(\linewidth - 22\tabcolsep) * \real{0.0833}}
  >{\raggedright\arraybackslash}p{(\linewidth - 22\tabcolsep) * \real{0.0833}}
  >{\raggedright\arraybackslash}p{(\linewidth - 22\tabcolsep) * \real{0.0833}}
  >{\raggedright\arraybackslash}p{(\linewidth - 22\tabcolsep) * \real{0.0833}}
  >{\raggedright\arraybackslash}p{(\linewidth - 22\tabcolsep) * \real{0.0833}}
  >{\raggedright\arraybackslash}p{(\linewidth - 22\tabcolsep) * \real{0.0833}}
  >{\raggedright\arraybackslash}p{(\linewidth - 22\tabcolsep) * \real{0.0833}}
  >{\raggedright\arraybackslash}p{(\linewidth - 22\tabcolsep) * \real{0.0833}}@{}}
\toprule\noalign{}
\begin{minipage}[b]{\linewidth}\raggedright
Jan
\end{minipage} & \begin{minipage}[b]{\linewidth}\raggedright
Fev
\end{minipage} & \begin{minipage}[b]{\linewidth}\raggedright
Mar
\end{minipage} & \begin{minipage}[b]{\linewidth}\raggedright
Abr
\end{minipage} & \begin{minipage}[b]{\linewidth}\raggedright
Mai
\end{minipage} & \begin{minipage}[b]{\linewidth}\raggedright
Jun
\end{minipage} & \begin{minipage}[b]{\linewidth}\raggedright
Jul
\end{minipage} & \begin{minipage}[b]{\linewidth}\raggedright
Ago
\end{minipage} & \begin{minipage}[b]{\linewidth}\raggedright
Set
\end{minipage} & \begin{minipage}[b]{\linewidth}\raggedright
Out
\end{minipage} & \begin{minipage}[b]{\linewidth}\raggedright
Nov
\end{minipage} & \begin{minipage}[b]{\linewidth}\raggedright
Dez
\end{minipage} \\
\midrule\noalign{}
\endhead
\bottomrule\noalign{}
\endlastfoot
6.56 & 6.11 & 5.31 & 4.26 & 3.15 & 2.69 & 2.99 & 3.69 & 4.36 & 5.27 &
6.39 & 6.73 \\
\end{longtable}
}

Média anual: 4.79 kWh/m²/dia. Inclinação recomendada: 27° Norte (anual);
37° inverno; 17° verão.

\textbf{Precipitação --- (mm/mês):}

{\def\LTcaptype{none} % do not increment counter
\begin{longtable}[]{@{}
  >{\raggedright\arraybackslash}p{(\linewidth - 22\tabcolsep) * \real{0.0833}}
  >{\raggedright\arraybackslash}p{(\linewidth - 22\tabcolsep) * \real{0.0833}}
  >{\raggedright\arraybackslash}p{(\linewidth - 22\tabcolsep) * \real{0.0833}}
  >{\raggedright\arraybackslash}p{(\linewidth - 22\tabcolsep) * \real{0.0833}}
  >{\raggedright\arraybackslash}p{(\linewidth - 22\tabcolsep) * \real{0.0833}}
  >{\raggedright\arraybackslash}p{(\linewidth - 22\tabcolsep) * \real{0.0833}}
  >{\raggedright\arraybackslash}p{(\linewidth - 22\tabcolsep) * \real{0.0833}}
  >{\raggedright\arraybackslash}p{(\linewidth - 22\tabcolsep) * \real{0.0833}}
  >{\raggedright\arraybackslash}p{(\linewidth - 22\tabcolsep) * \real{0.0833}}
  >{\raggedright\arraybackslash}p{(\linewidth - 22\tabcolsep) * \real{0.0833}}
  >{\raggedright\arraybackslash}p{(\linewidth - 22\tabcolsep) * \real{0.0833}}
  >{\raggedright\arraybackslash}p{(\linewidth - 22\tabcolsep) * \real{0.0833}}@{}}
\toprule\noalign{}
\begin{minipage}[b]{\linewidth}\raggedright
Jan
\end{minipage} & \begin{minipage}[b]{\linewidth}\raggedright
Fev
\end{minipage} & \begin{minipage}[b]{\linewidth}\raggedright
Mar
\end{minipage} & \begin{minipage}[b]{\linewidth}\raggedright
Abr
\end{minipage} & \begin{minipage}[b]{\linewidth}\raggedright
Mai
\end{minipage} & \begin{minipage}[b]{\linewidth}\raggedright
Jun
\end{minipage} & \begin{minipage}[b]{\linewidth}\raggedright
Jul
\end{minipage} & \begin{minipage}[b]{\linewidth}\raggedright
Ago
\end{minipage} & \begin{minipage}[b]{\linewidth}\raggedright
Set
\end{minipage} & \begin{minipage}[b]{\linewidth}\raggedright
Out
\end{minipage} & \begin{minipage}[b]{\linewidth}\raggedright
Nov
\end{minipage} & \begin{minipage}[b]{\linewidth}\raggedright
Dez
\end{minipage} \\
\midrule\noalign{}
\endhead
\bottomrule\noalign{}
\endlastfoot
156 & 125 & 149 & 186 & 147 & 112 & 92 & 90 & 122 & 215 & 190 & 188 \\
\end{longtable}
}

Média anual: \textasciitilde1.770 mm/ano. Estação chuvosa Out-Jan; seca
Jun-Ago.
\clearpage
\subsection*{Dados Departamentais}
\subsubsection{Desenvolvimento Humano e
Social}

\subsection{IDH Departamental}

{\def\LTcaptype{none} % do not increment counter
\begin{longtable}[]{@{}ll@{}}
\toprule\noalign{}
Indicador & Valor \\
\midrule\noalign{}
\endhead
\bottomrule\noalign{}
\endlastfoot
IDH & 0,728 \\
Classificação IDH & alto \\
Ranking nacional (1=melhor) & 4/18 \\
Esperança de vida & 74,0 anos \\
Anos de escolaridade média & 8,9 anos \\
RNB per capita (USD PPA) & 10.100 \\
\end{longtable}
}

\subsection{Indicadores de Pobreza}

{\def\LTcaptype{none} % do not increment counter
\begin{longtable}[]{@{}ll@{}}
\toprule\noalign{}
Indicador & Valor \\
\midrule\noalign{}
\endhead
\bottomrule\noalign{}
\endlastfoot
Taxa de pobreza total (\%) & 22,1\% \\
Taxa de pobreza extrema (\%) & 4,8\% \\
Índice de Gini & 0,432 \\
\end{longtable}
}

\subsection{Infraestrutura Básica}

{\def\LTcaptype{none} % do not increment counter
\begin{longtable}[]{@{}ll@{}}
\toprule\noalign{}
Indicador & Valor \\
\midrule\noalign{}
\endhead
\bottomrule\noalign{}
\endlastfoot
Acesso a água potável (\%) & 84,6\% \\
Acesso a saneamento (\%) & 83,1\% \\
\end{longtable}
}

\subsubsection{Segurança Pública}

\subsection{Indicadores}

{\def\LTcaptype{none} % do not increment counter
\begin{longtable}[]{@{}
  >{\raggedright\arraybackslash}p{(\linewidth - 2\tabcolsep) * \real{0.6111}}
  >{\raggedright\arraybackslash}p{(\linewidth - 2\tabcolsep) * \real{0.3889}}@{}}
\toprule\noalign{}
\begin{minipage}[b]{\linewidth}\raggedright
Indicador
\end{minipage} & \begin{minipage}[b]{\linewidth}\raggedright
Valor
\end{minipage} \\
\midrule\noalign{}
\endhead
\bottomrule\noalign{}
\endlastfoot
Taxa de homicídios (por 100k hab) & \textasciitilde4--6 (estimativa,
próxima à média nacional) \\
Crimes registrados (total/ano) & --- dados não desagregados
publicamente \\
Número de comisarías & \textasciitilde25 (estimativa --- departamento
populoso com fronteira) \\
Índice segurança relativo & médio \\
\end{longtable}
}

\subsubsection{Mercado de Terra Rural}

\subsection{Preços por Tipo de Terra}

{\def\LTcaptype{none} % do not increment counter
\begin{longtable}[]{@{}llll@{}}
\toprule\noalign{}
Tipo & Preço (USD/ha) & Faixa & Tendência \\
\midrule\noalign{}
\endhead
\bottomrule\noalign{}
\endlastfoot
Agrícola alta produtividade & 6.900 & 5.000 -- 10.000 & ↑ \\
Pastagem & 3.200 & 2.200 -- 4.500 & ↑ \\
Mata / reserva & 1.400 & 800 -- 2.200 & → \\
\end{longtable}
}

\subsection{Características do
Mercado}

{\def\LTcaptype{none} % do not increment counter
\begin{longtable}[]{@{}ll@{}}
\toprule\noalign{}
Parâmetro & Valor \\
\midrule\noalign{}
\endhead
\bottomrule\noalign{}
\endlastfoot
Valorização anual média (\%) & 7--9\% \\
Lote mínimo rural (ha) & 10 \\
Imposto territorial rural (\%) & 1\% sobre valor fiscal \\
Liquidez do mercado & alta \\
\end{longtable}
}
\clearpage
\newpage
\secaoDiagnostico{Alto Vera}{\mapaDistrito{0725}}{Alto Vera (Itapua) apresenta perfil operacional consolidado para
comparacao territorial, com leitura conjunta de infraestrutura, risco
hidroclimatico e capacidade de continuidade de servicos essenciais.

\subsubsection{Por que sim}

\begin{itemize}
\tightlist
\item
  Base institucional de referencia disponivel e rastreavel.
\item
  Estrutura comparativa (A-E/GSS) consistente para decisao relativa.
\item
  Plano de mitigacao acionavel ja definido no dossie.
\end{itemize}

\subsubsection{Por que não}

\begin{itemize}
\tightlist
\item
  Serie oficial distrital granular ainda pode apresentar lacunas
  temporais.
\item
  Sensibilidade do score exige monitoramento continuo de risco e
  conectividade.
\item
  Decisao final depende de validacao de campo e janela temporal recente.
\end{itemize}

\subsubsection{Pre-condicoes Minimas de
Decisao}

\begin{itemize}
\tightlist
\item
  Atualizar indicadores criticos em ciclo trimestral.
\item
  Confirmar trafegabilidade e contingencia hidrica/energetica local.
\item
  Validar checklist de resiliencia domiciliar/logistica antes de decisao
  final.
\end{itemize}

\subsubsection{Recomendação
Técnica}

Uso recomendado como candidato em analise multicriterio, condicionado ao
cumprimento das pre-condicoes acima. Referencia atual de score: GSS 8.0;
risco predominante: baixo.

\subsubsection{}

\begin{itemize}
\tightlist
\item
  \begin{enumerate}
  \def\labelenumi{\arabic{enumi})}
  \tightlist
  \item
    Delimitacao territorial validada por georreferencia institucional
    (Fonte: acesso: 2026-03-05).
  \end{enumerate}
\item
  \begin{enumerate}
  \def\labelenumi{\arabic{enumi})}
  \setcounter{enumi}{1}
  \tightlist
  \item
    Populacao municipal/distrital referenciada por base censitaria
    nacional (Fonte: acesso: 2026-03-05).
  \end{enumerate}
\item
  \begin{enumerate}
  \def\labelenumi{\arabic{enumi})}
  \setcounter{enumi}{2}
  \tightlist
  \item
    Comparabilidade intra-departamental apoiada em indicadores
    distritais oficiais (Fonte: acesso: 2026-03-05).
  \end{enumerate}
\item
  \begin{enumerate}
  \def\labelenumi{\arabic{enumi})}
  \setcounter{enumi}{3}
  \tightlist
  \item
    Estado macro de conectividade viaria ancorado em informacoes de
    rodovias (Fonte: acesso: 2026-03-05).
  \end{enumerate}
\item
  \begin{enumerate}
  \def\labelenumi{\arabic{enumi})}
  \setcounter{enumi}{4}
  \tightlist
  \item
    Exposicao hidrometeorologica monitorada por avisos oficiais (Fonte:
    acesso: 2026-03-05).
  \end{enumerate}
\item
  \begin{enumerate}
  \def\labelenumi{\arabic{enumi})}
  \setcounter{enumi}{5}
  \tightlist
  \item
    Historico de contexto meteorologico apoiado em publicacoes tecnicas
    (Fonte: acesso: 2026-03-05).
  \end{enumerate}
\item
  \begin{enumerate}
  \def\labelenumi{\arabic{enumi})}
  \setcounter{enumi}{6}
  \tightlist
  \item
    Capacidade institucional de resposta a eventos apoiada em acoes da
    SEN (Fonte: acesso: 2026-03-05).
  \end{enumerate}
\item
  \begin{enumerate}
  \def\labelenumi{\arabic{enumi})}
  \setcounter{enumi}{7}
  \tightlist
  \item
    Infraestrutura energetica analisada via operador nacional (Fonte:
    acesso: 2026-03-05).
  \end{enumerate}
\item
  \begin{enumerate}
  \def\labelenumi{\arabic{enumi})}
  \setcounter{enumi}{8}
  \tightlist
  \item
    Referencias de obras e logistica nacional verificadas no MOPC
    (Fonte: acesso: 2026-03-05).
  \end{enumerate}
\item
  \begin{enumerate}
  \def\labelenumi{\arabic{enumi})}
  \setcounter{enumi}{9}
  \tightlist
  \item
    Contexto sociopolitico institucional referenciado por autoridade
    eleitoral (Fonte: acesso: 2026-03-05).
  \end{enumerate}
\item
  \begin{enumerate}
  \def\labelenumi{\arabic{enumi})}
  \setcounter{enumi}{10}
  \tightlist
  \item
    Camada cartografica de apoio eleitoral/territorial consultada
    (Fonte: acesso: 2026-03-05).
  \end{enumerate}
\item
  \begin{enumerate}
  \def\labelenumi{\arabic{enumi})}
  \setcounter{enumi}{11}
  \tightlist
  \item
    Dados publicos complementares para verificacao cruzada (Fonte:
    acesso: 2026-03-05).
  \end{enumerate}
\end{itemize}}{dist:07_Itapua:Alto_Vera}
\begin{table}[ht!]
\small \begin{tabular}{p{3.0cm}p{0.8cm}p{6.0cm}} \toprule
\textbf{Critério} & \textbf{Nota} & \textbf{Justificativa} \\ \midrule
A - Ameaças Estratégicas & 8.8 & - \\
B - Risco Social & 7.5 & - \\
C - Riscos Naturais & 8.2 & - \\
D - Autossuficiência & 7.4 & - \\
E - Institucional & 8.0 & - \\
\midrule \textbf{GSS FINAL} & \textbf{8.0} & \textbf{Seguro (bom potencial com preparacao).} \\ \bottomrule \end{tabular}
\end{table}
\subsubsection{Vulnerabilidades e
Mitigação}\label{vuln-Alto_Vera-vulnerabilidades-e-mitigauxe7uxe3o}

\begin{itemize}
\tightlist
\item
  Exposicao hidrometeorologica controlada (\textless45/100).
\item
  Conectividade viaria favoravel (\textgreater=70/100).
\item
  Estoque de contingencia para 14 dias (agua/alimentos/combustivel),
  priorizado quando risco hidro=23/100.
\item
  Duplicar rotas/logistica quando conectividade viaria=84/100 for
  \textless70; manter rota secundaria mapeada e testada trimestralmente.
\item
  Plano de energia resiliente com autonomia minima de 72h
  (geracao/backup), revisao semestral.
\end{itemize}

Notas (0-10): - A: 8.8 - B: 7.5 - C: 8.2 - D: 7.4 - E: 8.0

Rastreabilidade das notas (base nos indicadores da secao 9): - A
(geoestrategia): combina conectividade viaria (84/100) e exposicao hidro
(23/100). - B (infraestrutura): combina conectividade (84/100) e escala
populacional (95027 hab). - C (riscos naturais): derivada da exposicao
hidrometeorologica (23/100). - D (autossuficiencia): combina
conectividade, ajuste de densidade (387.0 hab/km2) e risco hidro. - E
(ambiente sociopolitico): combinacao conservadora de conectividade,
escala populacional e risco hidro.

Calculo: - GSS = ((A x 2.5) + (B x 2.0) + (C x 1.5) + (D x 2.0) + (E x
2.0)) / 10 - GSS: 8.0

Classificacao: - Seguro (bom potencial com preparacao).

Sintese aplicada: - Dados textuais consolidados com base nas fontes
oficiais acima; proximos ciclos podem aprofundar indicadores numericos
especificos por distrito.
\subsubsection{Dossiê de Campo}\label{dossie-Alto_Vera-dossiuxea-de-campo}
\begin{description}[style=nextline,leftmargin=0.5cm,font=\bfseries]
\item[Coordenadas]  26°45′S 55°46′W (geocodificado via Nominatim).
\end{description}
\textbf{Fonte climática:} NASA POWER Climatology API (período 2001-2020)

\textbf{Fonte luminosa:} estimativa world atlas

\paragraph{Irradiação Solar
(kWh/m²/dia)}\label{irradiauxe7uxe3o-solar-kwhmuxb2dia}

{\def\LTcaptype{none} % do not increment counter
\begin{longtable}[]{@{}
  >{\raggedright\arraybackslash}p{(\linewidth - 24\tabcolsep) * \real{0.0746}}
  >{\raggedright\arraybackslash}p{(\linewidth - 24\tabcolsep) * \real{0.0746}}
  >{\raggedright\arraybackslash}p{(\linewidth - 24\tabcolsep) * \real{0.0746}}
  >{\raggedright\arraybackslash}p{(\linewidth - 24\tabcolsep) * \real{0.0746}}
  >{\raggedright\arraybackslash}p{(\linewidth - 24\tabcolsep) * \real{0.0746}}
  >{\raggedright\arraybackslash}p{(\linewidth - 24\tabcolsep) * \real{0.0746}}
  >{\raggedright\arraybackslash}p{(\linewidth - 24\tabcolsep) * \real{0.0746}}
  >{\raggedright\arraybackslash}p{(\linewidth - 24\tabcolsep) * \real{0.0746}}
  >{\raggedright\arraybackslash}p{(\linewidth - 24\tabcolsep) * \real{0.0746}}
  >{\raggedright\arraybackslash}p{(\linewidth - 24\tabcolsep) * \real{0.0746}}
  >{\raggedright\arraybackslash}p{(\linewidth - 24\tabcolsep) * \real{0.0746}}
  >{\raggedright\arraybackslash}p{(\linewidth - 24\tabcolsep) * \real{0.0746}}
  >{\raggedright\arraybackslash}p{(\linewidth - 24\tabcolsep) * \real{0.1045}}@{}}
\toprule\noalign{}
\begin{minipage}[b]{\linewidth}\raggedright
Jan
\end{minipage} & \begin{minipage}[b]{\linewidth}\raggedright
Fev
\end{minipage} & \begin{minipage}[b]{\linewidth}\raggedright
Mar
\end{minipage} & \begin{minipage}[b]{\linewidth}\raggedright
Abr
\end{minipage} & \begin{minipage}[b]{\linewidth}\raggedright
Mai
\end{minipage} & \begin{minipage}[b]{\linewidth}\raggedright
Jun
\end{minipage} & \begin{minipage}[b]{\linewidth}\raggedright
Jul
\end{minipage} & \begin{minipage}[b]{\linewidth}\raggedright
Ago
\end{minipage} & \begin{minipage}[b]{\linewidth}\raggedright
Set
\end{minipage} & \begin{minipage}[b]{\linewidth}\raggedright
Out
\end{minipage} & \begin{minipage}[b]{\linewidth}\raggedright
Nov
\end{minipage} & \begin{minipage}[b]{\linewidth}\raggedright
Dez
\end{minipage} & \begin{minipage}[b]{\linewidth}\raggedright
Média
\end{minipage} \\
\midrule\noalign{}
\endhead
\bottomrule\noalign{}
\endlastfoot
6.56 & 6.06 & 5.32 & 4.33 & 3.21 & 2.76 & 3.08 & 3.81 & 4.39 & 5.25 &
6.31 & 6.67 & \textbf{4.81} \\
\end{longtable}
}

\textbf{Inclinação solar recomendada:} 27° N (anual) · 37° N (inverno
jun-ago) · 17° N (verão nov-jan)

\paragraph{Precipitação (mm/mês)}\label{precipitauxe7uxe3o-mmmuxeas}

{\def\LTcaptype{none} % do not increment counter
\begin{longtable}[]{@{}
  >{\raggedright\arraybackslash}p{(\linewidth - 24\tabcolsep) * \real{0.0704}}
  >{\raggedright\arraybackslash}p{(\linewidth - 24\tabcolsep) * \real{0.0704}}
  >{\raggedright\arraybackslash}p{(\linewidth - 24\tabcolsep) * \real{0.0704}}
  >{\raggedright\arraybackslash}p{(\linewidth - 24\tabcolsep) * \real{0.0704}}
  >{\raggedright\arraybackslash}p{(\linewidth - 24\tabcolsep) * \real{0.0704}}
  >{\raggedright\arraybackslash}p{(\linewidth - 24\tabcolsep) * \real{0.0704}}
  >{\raggedright\arraybackslash}p{(\linewidth - 24\tabcolsep) * \real{0.0704}}
  >{\raggedright\arraybackslash}p{(\linewidth - 24\tabcolsep) * \real{0.0704}}
  >{\raggedright\arraybackslash}p{(\linewidth - 24\tabcolsep) * \real{0.0704}}
  >{\raggedright\arraybackslash}p{(\linewidth - 24\tabcolsep) * \real{0.0704}}
  >{\raggedright\arraybackslash}p{(\linewidth - 24\tabcolsep) * \real{0.0704}}
  >{\raggedright\arraybackslash}p{(\linewidth - 24\tabcolsep) * \real{0.0704}}
  >{\raggedright\arraybackslash}p{(\linewidth - 24\tabcolsep) * \real{0.1549}}@{}}
\toprule\noalign{}
\begin{minipage}[b]{\linewidth}\raggedright
Jan
\end{minipage} & \begin{minipage}[b]{\linewidth}\raggedright
Fev
\end{minipage} & \begin{minipage}[b]{\linewidth}\raggedright
Mar
\end{minipage} & \begin{minipage}[b]{\linewidth}\raggedright
Abr
\end{minipage} & \begin{minipage}[b]{\linewidth}\raggedright
Mai
\end{minipage} & \begin{minipage}[b]{\linewidth}\raggedright
Jun
\end{minipage} & \begin{minipage}[b]{\linewidth}\raggedright
Jul
\end{minipage} & \begin{minipage}[b]{\linewidth}\raggedright
Ago
\end{minipage} & \begin{minipage}[b]{\linewidth}\raggedright
Set
\end{minipage} & \begin{minipage}[b]{\linewidth}\raggedright
Out
\end{minipage} & \begin{minipage}[b]{\linewidth}\raggedright
Nov
\end{minipage} & \begin{minipage}[b]{\linewidth}\raggedright
Dez
\end{minipage} & \begin{minipage}[b]{\linewidth}\raggedright
Total/ano
\end{minipage} \\
\midrule\noalign{}
\endhead
\bottomrule\noalign{}
\endlastfoot
154 & 122 & 146 & 184 & 150 & 110 & 90 & 82 & 113 & 206 & 193 & 180 &
\textbf{1732 mm} \\
\end{longtable}
}

\paragraph{Poluição Luminosa}\label{poluiuxe7uxe3o-luminosa}

{\def\LTcaptype{none} % do not increment counter
\begin{longtable}[]{@{}ll@{}}
\toprule\noalign{}
Parâmetro & Valor \\
\midrule\noalign{}
\endhead
\bottomrule\noalign{}
\endlastfoot
Escala Bortle & 5 --- Céu suburbano \\
Radiância artificial & 5.0 nW/cm²/sr \\
\end{longtable}
}
\paragraph{Solo (SoilGrids 2.0, média ponderada 0--30
cm)}

{\def\LTcaptype{none} % do not increment counter
\begin{longtable}[]{@{}ll@{}}
\toprule\noalign{}
Parâmetro & Valor \\
\midrule\noalign{}
\endhead
\bottomrule\noalign{}
\endlastfoot
pH (H₂O) & 5.2 \\
Carbono orgânico (SOC) & 22.42 g/kg \\
Argila & 42.58 \% \\
Areia & 24.54 \% \\
Silte (calc.) & 32.9 \% \\
Densidade aparente & 1.23 g/cm³ \\
\textbf{Aptidão agrícola} & \textbf{Média} \\
\end{longtable}
}

Fonte: ISRIC SoilGrids 2.0 via WCS. Coords: -26.75°, -55.7667°. Média
ponderada camadas 0-5, 5-15, 15-30 cm.
\subsubsection{Indicadores Sociais}

\textbf{Fontes:} PNUD 2020, INE Censo 2022, INE EPHC 2023, Ministerio
Público 2024\\
\textbf{Nota:} Valores marcados como \emph{(dept.)} referem-se ao
departamento de Itapúa; valores \emph{(dist.)} são específicos deste
distrito. Dados marcados \emph{(est.)} são estimativas.

{\def\LTcaptype{none} % do not increment counter
\begin{longtable}[]{@{}
  >{\raggedright\arraybackslash}p{(\linewidth - 4\tabcolsep) * \real{0.4231}}
  >{\raggedright\arraybackslash}p{(\linewidth - 4\tabcolsep) * \real{0.2692}}
  >{\raggedright\arraybackslash}p{(\linewidth - 4\tabcolsep) * \real{0.3077}}@{}}
\toprule\noalign{}
\begin{minipage}[b]{\linewidth}\raggedright
Indicador
\end{minipage} & \begin{minipage}[b]{\linewidth}\raggedright
Valor
\end{minipage} & \begin{minipage}[b]{\linewidth}\raggedright
Âmbito
\end{minipage} \\
\midrule\noalign{}
\endhead
\bottomrule\noalign{}
\endlastfoot
IDH (2020) & 0,728 & dept. \\
Ranking IDH nacional & 4/18 & dept. \\
Esperança de vida & 74,0 anos & dept. \\
Escolaridade média & 6.3 anos & dist. \\
RNB per capita (USD PPA) & 10.100 & dept. \\
Pobreza monetária (\%) & 22,1\% & dept. \\
Pobreza extrema (\%) & 4,8\% & dept. \\
Índice de Gini & 0,432 & dept. \\
Acesso a água potável (\%) & 84,6\% & dept. \\
Acesso a saneamento (\%) & 83,1\% & dept. \\
Taxa de homicídios (est., /100k hab) & \textasciitilde4--6 (estimativa,
próxima à média nacional) & dept. \\
Índice de segurança & médio & dept. \\
População (Censo 2022) & 11.503 hab. & dist. \\
Idade mediana & 25.0 anos & dist. \\
Taxa de fecundidade & N/D & dist. \\
\end{longtable}
}
\subsubsection{Combustível}

\textbf{Referência:} PETROPAR / postos locais (2024) \textbf{Tipo de
localidade:} Interior

{\def\LTcaptype{none} % do not increment counter
\begin{longtable}[]{@{}lll@{}}
\toprule\noalign{}
Combustível & USD/litro & Gs/litro (aprox.) \\
\midrule\noalign{}
\endhead
\bottomrule\noalign{}
\endlastfoot
Gasolina 93 oct & 0.98 & 7,252 \\
Gasolina 97 oct (premium) & 1.08 & 7,992 \\
Diesel & 0.91 & 6,734 \\
\end{longtable}
}

\begin{quote}
Preços podem variar ±5\% conforme posto e sazonalidade. Chaco e interior
remoto apresentam maior variação.
\end{quote}

\subsubsection{Cobertura Celular}

\textbf{Fonte:} CONATEL PY / operadoras (2024)

{\def\LTcaptype{none} % do not increment counter
\begin{longtable}[]{@{}ll@{}}
\toprule\noalign{}
Parâmetro & Valor \\
\midrule\noalign{}
\endhead
\bottomrule\noalign{}
\endlastfoot
Cobertura 4G população (dept.) & 92\% \\
Cobertura 4G área rural & 82\% \\
Melhor operadora & Tigo \\
Qualidade rural & boa \\
\end{longtable}
}

\begin{quote}
Para áreas rurais fora do núcleo urbano, recomenda-se chip Tigo como
principal e Personal como backup.
\end{quote}

\subsubsection{Internet}

\textbf{Fonte:} CONATEL / Speedtest Ookla (2024)

{\def\LTcaptype{none} % do not increment counter
\begin{longtable}[]{@{}ll@{}}
\toprule\noalign{}
Parâmetro & Valor \\
\midrule\noalign{}
\endhead
\bottomrule\noalign{}
\endlastfoot
Velocidade média download & 75 Mbps \\
Domicílios com internet (dept.) & 78\% \\
Tecnologia predominante & fibra/rádio \\
Opção rural & Starlink disponível (\textasciitilde USD 44/mês) \\
\end{longtable}
}

\subsubsection{Mercado Imobiliário e Terra
Rural}

\textbf{Fonte:} INDERT / Clasificados.com.py (2024)

{\def\LTcaptype{none} % do not increment counter
\begin{longtable}[]{@{}ll@{}}
\toprule\noalign{}
Tipo & Referência \\
\midrule\noalign{}
\endhead
\bottomrule\noalign{}
\endlastfoot
Terra agrícola alta prod. (USD/ha) & 6,900 \\
Imóvel urbano (USD/m²) & 1,300 \\
Aluguel 2 quartos (USD/mês) & 550 \\
\end{longtable}
}

\begin{quote}
Valores de referência departamental. Localidades menores podem ter
preços 20--40\% abaixo da capital departamental.
\end{quote}

\subsubsection{Saúde}

\textbf{Fonte:} MSPBS / IPS Paraguay (2024-2026), consolidação
departamental e proxy local

{\def\LTcaptype{none} % do not increment counter
\begin{longtable}[]{@{}ll@{}}
\toprule\noalign{}
Serviço & Disponibilidade \\
\midrule\noalign{}
\endhead
\bottomrule\noalign{}
\endlastfoot
USF / Posto de Saúde & sim \\
Hospital Regional & não \\
IPS (seguro social) & sim \\
Farmácia & sim \\
Distância ao hospital de referência & \textasciitilde82 km
(Encarnacion) \\
\end{longtable}
}

\textbf{Principais estabelecimentos:} USF local; referencia hospitalar
em Encarnacion

\textbf{Observação para imigrantes:} Atendimento primario local ou em
raio curto; casos de maior complexidade seguem para Encarnacion.
Cobertura privada continua recomendada para especialidades e urgências
de maior complexidade.
\newpage
\secaoDiagnostico{Bella Vista}{\mapaDistrito{0702}}{Bella Vista é um dos pilares das Colonias Unidas. Oferece uma qualidade
de vida e segurança social excepcionais, apoiadas em um sistema
cooperativista robusto e solos de altíssima produtividade.

\subsubsection{Por que sim}

\begin{itemize}
\tightlist
\item
  Sistema cooperativista exemplar (segurança econômica).
\item
  Solo de altíssima fertilidade (Terra Roxa).
\item
  Elevada coesão social.
\end{itemize}

\subsubsection{Por que não}

\begin{itemize}
\tightlist
\item
  Exposição à fronteira e a rotas de tráfego intenso.
\end{itemize}

\subsubsection{Recomendação
Técnica}

Altamente recomendado para quem busca integração comunitária e produção
agroindustrial em ambiente seguro.

\subsubsection{Analise de Sensibilidade}

\begin{itemize}
\tightlist
\item
  Cenario otimista: GSS 6.4 (Expansão da agroindústria local).
\item
  Cenario conservador: GSS 5.4 (Crise nos preços internacionais de
  commodities).
\end{itemize}

\subsubsection{}

\begin{itemize}
\tightlist
\item
  Dados populacionais INE 2022.
\item
  Registros de produção das cooperativas locais.
\item
  Dados de infraestrutura viária MOPC.
\end{itemize}}{dist:07_Itapua:Bella_Vista}
\begin{table}[ht!]
\small \begin{tabular}{p{3.0cm}p{0.8cm}p{6.0cm}} \toprule
\textbf{Critério} & \textbf{Nota} & \textbf{Justificativa} \\ \midrule
A - Ameaças Estratégicas & 5.0 & - \\
B - Risco Social & 6.0 & - \\
C - Riscos Naturais & 7.0 & - \\
D - Autossuficiência & 6.0 & - \\
E - Institucional & 6.0 & - \\
\midrule \textbf{GSS FINAL} & \textbf{5.9} & \textbf{Seguro (Potencial moderado de resiliência).} \\ \bottomrule \end{tabular}
\end{table}
\subsubsection{Vulnerabilidades e
Mitigação}\label{vuln-Bella_Vista-vulnerabilidades-e-mitigauxe7uxe3o}

\begin{itemize}
\item
  \textbf{Fronteira:} Exposição a fluxos transfronteiriços. Mitigação:
  Monitoramento e integração com forças de segurança.
\item
  \textbf{Monoculturas:} Dependência econômica da erva-mate e soja.
  Mitigação: Diversificação de culturas na propriedade.
\item
  INE Censo 2022: 10.447 habitantes.
\item
  Solo: Terra Roxa de origem basáltica.
\item
  Homicídios: \textless5\% do total nacional (Itapúa).
\end{itemize}
\subsubsection{Dossiê de Campo}\label{dossie-Bella_Vista-dossiuxea-de-campo}
\begin{description}[style=nextline,leftmargin=0.5cm,font=\bfseries]
\item[Coordenadas]  27°02′S 55°34′W.
\item[Topografia]  Relevo ondulado, margens do Rio Paraná. Área de \textasciitilde 285 km².
\item[Alvos Estrategicos]  Centro de produção de erva-mate e agroindústria. Proximidade com a fronteira argentina. Sem alvos militares diretos.
\item[Fallout]  Ventos predominantes do Norte.
\item[Populacao]  10.447 (Censo 2022 Final).
\item[Densidade]  \textasciitilde 36,7 hab/km².
\item[Vias de Acesso]  Localizado na Ruta PY06, principal eixo Itapúa-Alto Paraná.
\item[Servicos]  Excelente infraestrutura de serviços privados e cooperativismo forte (Colonias Unidas).
\item[Hidrologia]  Proximidade com o Rio Paraná exige atenção a áreas ribeirinhas, mas a zona urbana é elevada.
\item[Sismicidade]  Baixa.
\item[Solo]  Terra Roxa (Oxisols) de altíssima fertilidade.
\item[Agua]  Sistemas de cooperativas eficientes.
\item[Energia]  Rede ANDE estável.
\item[Producao]  Erva-mate, soja, trigo e agroindústria de alimentos.
\item[Seguranca]  Região das Colonias Unidas, conhecida por altos níveis de segurança e ordem social.
\item[Clima Social]  Forte influência de imigrantes europeus, cultura de trabalho e cooperativismo.
\end{description}
\textbf{Fonte climática:} NASA POWER Climatology API (período 2001-2020)

\textbf{Fonte luminosa:} estimativa world atlas

\paragraph{Irradiação Solar
(kWh/m²/dia)}\label{irradiauxe7uxe3o-solar-kwhmuxb2dia}

{\def\LTcaptype{none} % do not increment counter
\begin{longtable}[]{@{}
  >{\raggedright\arraybackslash}p{(\linewidth - 24\tabcolsep) * \real{0.0746}}
  >{\raggedright\arraybackslash}p{(\linewidth - 24\tabcolsep) * \real{0.0746}}
  >{\raggedright\arraybackslash}p{(\linewidth - 24\tabcolsep) * \real{0.0746}}
  >{\raggedright\arraybackslash}p{(\linewidth - 24\tabcolsep) * \real{0.0746}}
  >{\raggedright\arraybackslash}p{(\linewidth - 24\tabcolsep) * \real{0.0746}}
  >{\raggedright\arraybackslash}p{(\linewidth - 24\tabcolsep) * \real{0.0746}}
  >{\raggedright\arraybackslash}p{(\linewidth - 24\tabcolsep) * \real{0.0746}}
  >{\raggedright\arraybackslash}p{(\linewidth - 24\tabcolsep) * \real{0.0746}}
  >{\raggedright\arraybackslash}p{(\linewidth - 24\tabcolsep) * \real{0.0746}}
  >{\raggedright\arraybackslash}p{(\linewidth - 24\tabcolsep) * \real{0.0746}}
  >{\raggedright\arraybackslash}p{(\linewidth - 24\tabcolsep) * \real{0.0746}}
  >{\raggedright\arraybackslash}p{(\linewidth - 24\tabcolsep) * \real{0.0746}}
  >{\raggedright\arraybackslash}p{(\linewidth - 24\tabcolsep) * \real{0.1045}}@{}}
\toprule\noalign{}
\begin{minipage}[b]{\linewidth}\raggedright
Jan
\end{minipage} & \begin{minipage}[b]{\linewidth}\raggedright
Fev
\end{minipage} & \begin{minipage}[b]{\linewidth}\raggedright
Mar
\end{minipage} & \begin{minipage}[b]{\linewidth}\raggedright
Abr
\end{minipage} & \begin{minipage}[b]{\linewidth}\raggedright
Mai
\end{minipage} & \begin{minipage}[b]{\linewidth}\raggedright
Jun
\end{minipage} & \begin{minipage}[b]{\linewidth}\raggedright
Jul
\end{minipage} & \begin{minipage}[b]{\linewidth}\raggedright
Ago
\end{minipage} & \begin{minipage}[b]{\linewidth}\raggedright
Set
\end{minipage} & \begin{minipage}[b]{\linewidth}\raggedright
Out
\end{minipage} & \begin{minipage}[b]{\linewidth}\raggedright
Nov
\end{minipage} & \begin{minipage}[b]{\linewidth}\raggedright
Dez
\end{minipage} & \begin{minipage}[b]{\linewidth}\raggedright
Média
\end{minipage} \\
\midrule\noalign{}
\endhead
\bottomrule\noalign{}
\endlastfoot
6.56 & 6.11 & 5.31 & 4.26 & 3.15 & 2.69 & 2.99 & 3.69 & 4.36 & 5.27 &
6.39 & 6.73 & \textbf{4.79} \\
\end{longtable}
}

\textbf{Inclinação solar recomendada:} 27° N (anual) · 37° N (inverno
jun-ago) · 17° N (verão nov-jan)

\paragraph{Precipitação (mm/mês)}\label{precipitauxe7uxe3o-mmmuxeas}

{\def\LTcaptype{none} % do not increment counter
\begin{longtable}[]{@{}
  >{\raggedright\arraybackslash}p{(\linewidth - 24\tabcolsep) * \real{0.0704}}
  >{\raggedright\arraybackslash}p{(\linewidth - 24\tabcolsep) * \real{0.0704}}
  >{\raggedright\arraybackslash}p{(\linewidth - 24\tabcolsep) * \real{0.0704}}
  >{\raggedright\arraybackslash}p{(\linewidth - 24\tabcolsep) * \real{0.0704}}
  >{\raggedright\arraybackslash}p{(\linewidth - 24\tabcolsep) * \real{0.0704}}
  >{\raggedright\arraybackslash}p{(\linewidth - 24\tabcolsep) * \real{0.0704}}
  >{\raggedright\arraybackslash}p{(\linewidth - 24\tabcolsep) * \real{0.0704}}
  >{\raggedright\arraybackslash}p{(\linewidth - 24\tabcolsep) * \real{0.0704}}
  >{\raggedright\arraybackslash}p{(\linewidth - 24\tabcolsep) * \real{0.0704}}
  >{\raggedright\arraybackslash}p{(\linewidth - 24\tabcolsep) * \real{0.0704}}
  >{\raggedright\arraybackslash}p{(\linewidth - 24\tabcolsep) * \real{0.0704}}
  >{\raggedright\arraybackslash}p{(\linewidth - 24\tabcolsep) * \real{0.0704}}
  >{\raggedright\arraybackslash}p{(\linewidth - 24\tabcolsep) * \real{0.1549}}@{}}
\toprule\noalign{}
\begin{minipage}[b]{\linewidth}\raggedright
Jan
\end{minipage} & \begin{minipage}[b]{\linewidth}\raggedright
Fev
\end{minipage} & \begin{minipage}[b]{\linewidth}\raggedright
Mar
\end{minipage} & \begin{minipage}[b]{\linewidth}\raggedright
Abr
\end{minipage} & \begin{minipage}[b]{\linewidth}\raggedright
Mai
\end{minipage} & \begin{minipage}[b]{\linewidth}\raggedright
Jun
\end{minipage} & \begin{minipage}[b]{\linewidth}\raggedright
Jul
\end{minipage} & \begin{minipage}[b]{\linewidth}\raggedright
Ago
\end{minipage} & \begin{minipage}[b]{\linewidth}\raggedright
Set
\end{minipage} & \begin{minipage}[b]{\linewidth}\raggedright
Out
\end{minipage} & \begin{minipage}[b]{\linewidth}\raggedright
Nov
\end{minipage} & \begin{minipage}[b]{\linewidth}\raggedright
Dez
\end{minipage} & \begin{minipage}[b]{\linewidth}\raggedright
Total/ano
\end{minipage} \\
\midrule\noalign{}
\endhead
\bottomrule\noalign{}
\endlastfoot
154 & 122 & 146 & 184 & 150 & 110 & 90 & 82 & 113 & 206 & 193 & 180 &
\textbf{1732 mm} \\
\end{longtable}
}

\paragraph{Poluição Luminosa}\label{poluiuxe7uxe3o-luminosa}

{\def\LTcaptype{none} % do not increment counter
\begin{longtable}[]{@{}ll@{}}
\toprule\noalign{}
Parâmetro & Valor \\
\midrule\noalign{}
\endhead
\bottomrule\noalign{}
\endlastfoot
Escala Bortle & 5 --- Céu suburbano \\
Radiância artificial & 5.0 nW/cm²/sr \\
\end{longtable}
}
\paragraph{Solo (SoilGrids 2.0, média ponderada 0--30
cm)}

{\def\LTcaptype{none} % do not increment counter
\begin{longtable}[]{@{}ll@{}}
\toprule\noalign{}
Parâmetro & Valor \\
\midrule\noalign{}
\endhead
\bottomrule\noalign{}
\endlastfoot
pH (H₂O) & 5.2 \\
Carbono orgânico (SOC) & 20.73 g/kg \\
Argila & 37.7 \% \\
Areia & 23.73 \% \\
Silte (calc.) & 38.6 \% \\
Densidade aparente & 1.3 g/cm³ \\
\textbf{Aptidão agrícola} & \textbf{Média} \\
\end{longtable}
}

Fonte: ISRIC SoilGrids 2.0 via WCS. Coords: -27.0333°, -55.5667°. Média
ponderada camadas 0-5, 5-15, 15-30 cm.

\begin{itemize}
\tightlist
\item
  \textbf{Agua:} Sistemas de cooperativas eficientes.
\item
  \textbf{Energia:} Rede ANDE estável.
\item
  \textbf{Producao:} Erva-mate, soja, trigo e agroindústria de
  alimentos.
\end{itemize}
\subsubsection{Indicadores Sociais}

\textbf{Fontes:} PNUD 2020, INE Censo 2022, INE EPHC 2023, Ministerio
Público 2024\\
\textbf{Nota:} Valores marcados como \emph{(dept.)} referem-se ao
departamento de Itapúa; valores \emph{(dist.)} são específicos deste
distrito. Dados marcados \emph{(est.)} são estimativas.

{\def\LTcaptype{none} % do not increment counter
\begin{longtable}[]{@{}
  >{\raggedright\arraybackslash}p{(\linewidth - 4\tabcolsep) * \real{0.4231}}
  >{\raggedright\arraybackslash}p{(\linewidth - 4\tabcolsep) * \real{0.2692}}
  >{\raggedright\arraybackslash}p{(\linewidth - 4\tabcolsep) * \real{0.3077}}@{}}
\toprule\noalign{}
\begin{minipage}[b]{\linewidth}\raggedright
Indicador
\end{minipage} & \begin{minipage}[b]{\linewidth}\raggedright
Valor
\end{minipage} & \begin{minipage}[b]{\linewidth}\raggedright
Âmbito
\end{minipage} \\
\midrule\noalign{}
\endhead
\bottomrule\noalign{}
\endlastfoot
IDH (2020) & 0,728 & dept. \\
Ranking IDH nacional & 4/18 & dept. \\
Esperança de vida & 74,0 anos & dept. \\
Escolaridade média & 7.0 anos & dist. \\
RNB per capita (USD PPA) & 10.100 & dept. \\
Pobreza monetária (\%) & 22,1\% & dept. \\
Pobreza extrema (\%) & 4,8\% & dept. \\
Índice de Gini & 0,432 & dept. \\
Acesso a água potável (\%) & 84,6\% & dept. \\
Acesso a saneamento (\%) & 83,1\% & dept. \\
Taxa de homicídios (est., /100k hab) & \textasciitilde4--6 (estimativa,
próxima à média nacional) & dept. \\
Índice de segurança & médio & dept. \\
População (Censo 2022) & 10.447 hab. & dist. \\
Idade mediana & 29.0 anos & dist. \\
Taxa de fecundidade & N/D & dist. \\
\end{longtable}
}
\subsubsection{Combustível}

\textbf{Referência:} PETROPAR / postos locais (2024) \textbf{Tipo de
localidade:} Interior

{\def\LTcaptype{none} % do not increment counter
\begin{longtable}[]{@{}lll@{}}
\toprule\noalign{}
Combustível & USD/litro & Gs/litro (aprox.) \\
\midrule\noalign{}
\endhead
\bottomrule\noalign{}
\endlastfoot
Gasolina 93 oct & 0.98 & 7,252 \\
Gasolina 97 oct (premium) & 1.08 & 7,992 \\
Diesel & 0.91 & 6,734 \\
\end{longtable}
}

\begin{quote}
Preços podem variar ±5\% conforme posto e sazonalidade. Chaco e interior
remoto apresentam maior variação.
\end{quote}

\subsubsection{Cobertura Celular}

\textbf{Fonte:} CONATEL PY / operadoras (2024)

{\def\LTcaptype{none} % do not increment counter
\begin{longtable}[]{@{}ll@{}}
\toprule\noalign{}
Parâmetro & Valor \\
\midrule\noalign{}
\endhead
\bottomrule\noalign{}
\endlastfoot
Cobertura 4G população (dept.) & 92\% \\
Cobertura 4G área rural & 82\% \\
Melhor operadora & Tigo \\
Qualidade rural & boa \\
\end{longtable}
}

\begin{quote}
Para áreas rurais fora do núcleo urbano, recomenda-se chip Tigo como
principal e Personal como backup.
\end{quote}

\subsubsection{Internet}

\textbf{Fonte:} CONATEL / Speedtest Ookla (2024)

{\def\LTcaptype{none} % do not increment counter
\begin{longtable}[]{@{}ll@{}}
\toprule\noalign{}
Parâmetro & Valor \\
\midrule\noalign{}
\endhead
\bottomrule\noalign{}
\endlastfoot
Velocidade média download & 75 Mbps \\
Domicílios com internet (dept.) & 78\% \\
Tecnologia predominante & fibra/rádio \\
Opção rural & Starlink disponível (\textasciitilde USD 44/mês) \\
\end{longtable}
}

\subsubsection{Mercado Imobiliário e Terra
Rural}

\textbf{Fonte:} INDERT / Clasificados.com.py (2024)

{\def\LTcaptype{none} % do not increment counter
\begin{longtable}[]{@{}ll@{}}
\toprule\noalign{}
Tipo & Referência \\
\midrule\noalign{}
\endhead
\bottomrule\noalign{}
\endlastfoot
Terra agrícola alta prod. (USD/ha) & 6,900 \\
Imóvel urbano (USD/m²) & 1,300 \\
Aluguel 2 quartos (USD/mês) & 550 \\
\end{longtable}
}

\begin{quote}
Valores de referência departamental. Localidades menores podem ter
preços 20--40\% abaixo da capital departamental.
\end{quote}

\subsubsection{Saúde}

\textbf{Fonte:} MSPBS / IPS Paraguay (2024-2026), consolidação
departamental e proxy local

{\def\LTcaptype{none} % do not increment counter
\begin{longtable}[]{@{}ll@{}}
\toprule\noalign{}
Serviço & Disponibilidade \\
\midrule\noalign{}
\endhead
\bottomrule\noalign{}
\endlastfoot
USF / Posto de Saúde & sim \\
Hospital Regional & não \\
IPS (seguro social) & sim \\
Farmácia & sim \\
Distância ao hospital de referência & \textasciitilde55 km
(Encarnacion) \\
\end{longtable}
}

\textbf{Principais estabelecimentos:} USF local; referencia hospitalar
em Encarnacion

\textbf{Observação para imigrantes:} Atendimento primario local ou em
raio curto; casos de maior complexidade seguem para Encarnacion.
Cobertura privada continua recomendada para especialidades e urgências
de maior complexidade.
\newpage
\secaoDiagnostico{Cambreta}{\mapaConteudo{-27.3333}{-55.8333}}{Cambyretá funciona como uma extensão urbana estratégica de Encarnación.
Oferece as melhores facilidades em termos de serviços e conectividade,
mas sua alta densidade populacional e proximidade com o centro regional
reduzem a resiliência em cenários de colapso urbano sistêmico.

\subsubsection{Por que sim}

\begin{itemize}
\tightlist
\item
  Conectividade excepcional e acesso a serviços de alta complexidade.
\item
  Solo fértil para horticultura de apoio.
\end{itemize}

\subsubsection{Por que não}

\begin{itemize}
\tightlist
\item
  Alta densidade demográfica (251 hab/km²).
\item
  Dependência excessiva da infraestrutura de Encarnación.
\end{itemize}

\subsubsection{Recomendação
Técnica}

Ideal para quem busca as conveniências da vida urbana em uma região
segura, desde que possua plano de contingência para crises de
abastecimento.

\subsubsection{Analise de Sensibilidade}

\begin{itemize}
\tightlist
\item
  Cenario otimista: GSS 6.3 (Melhoria na drenagem urbana e
  infraestrutura local própria).
\item
  Cenario conservador: GSS 5.4 (Insegurança urbana crescente por
  transbordamento da metrópole).
\end{itemize}

\subsubsection{}

\begin{itemize}
\tightlist
\item
  Dados populacionais INE 2022.
\item
  Mapeamento urbano e de serviços municipais.
\item
  Dados de rodovias MOPC.
\end{itemize}}{dist:07_Itapua:Cambreta}
\begin{table}[ht!]
\small \begin{tabular}{p{3.0cm}p{0.8cm}p{6.0cm}} \toprule
\textbf{Critério} & \textbf{Nota} & \textbf{Justificativa} \\ \midrule
A - Ameaças Estratégicas & 5.0 & - \\
B - Risco Social & 7.0 & - \\
C - Riscos Naturais & 7.0 & - \\
D - Autossuficiência & 5.0 & - \\
E - Institucional & 6.0 & - \\
\midrule \textbf{GSS FINAL} & \textbf{5.9} & \textbf{Seguro (Potencial moderado de resiliência).} \\ \bottomrule \end{tabular}
\end{table}
\subsubsection{Vulnerabilidades e
Mitigação}\label{vuln-Cambreta-vulnerabilidades-e-mitigauxe7uxe3o}

\begin{itemize}
\item
  \textbf{Densidade Populacional:} Risco de saturação de serviços
  básicos em crises migratórias ou sanitárias. Mitigação: Estoque de
  emergência individual e plano de evacuação para áreas rurais.
\item
  \textbf{Dependência Urbana:} Elevada dependência de Encarnación para
  infraestrutura crítica. Mitigação: Mapear rotas de fuga alternativas
  evitando os eixos centrais.
\item
  INE Censo 2022: 47.717 habitantes.
\item
  Solo: Terra Roxa de alta fertilidade.
\item
  Homicídios: Itapúa possui uma das menores taxas nacionais.
\end{itemize}
\subsubsection{Dossiê de Campo}\label{dossie-Cambreta-dossiuxea-de-campo}
\begin{description}[style=nextline,leftmargin=0.5cm,font=\bfseries]
\item[Coordenadas]  27°20′S 55°50′W.
\item[Topografia]  Relevo ondulado com áreas de várzea nas proximidades de arroios tributários do Rio Paraná. Área de \textasciitilde 190 km².
\item[Alvos Estrategicos]  Distrito limítrofe com Encarnación. Ponto de expansão urbana e logística. Sem alvos militares.
\item[Fallout]  Ventos predominantes do Norte/Nordeste.
\item[Populacao]  47.717 (Censo 2022 Final). Segundo distrito mais populoso de Itapúa.
\item[Densidade]  \textasciitilde 251,1 hab/km².
\item[Vias de Acesso]  Excelente conectividade viária, integrado à malha urbana de Encarnación.
\item[Servicos]  Acesso direto a todos os serviços da capital regional (Encarnación).
\item[Hidrologia]  Diversos arroios. Risco localizado de alagamentos em áreas de ocupação irregular ou baixadas em chuvas intensas.
\item[Sismicidade]  Baixa.
\item[Solo]  Terra Roxa de origem basáltica, altamente fértil, mas sob forte pressão de urbanização.
\item[Agua]  Cobertura de rede de água e poços comunitários.
\item[Energia]  Rede ANDE estável.
\item[Producao]  Olericultura, pequenas granjas e serviços urbanos.
\item[Seguranca]  Perfil urbano-suburbano. Segurança estável, mas requer atenção a índices de criminalidade urbana típicos de áreas metropolitanas.
\item[Clima Social]  Sociedade diversa, integração total com Encarnación.
\end{description}
\textbf{Fonte climática:} NASA POWER Climatology API (período 2001-2020)

\textbf{Fonte luminosa:} estimativa world atlas

\paragraph{Irradiação Solar
(kWh/m²/dia)}\label{irradiauxe7uxe3o-solar-kwhmuxb2dia}

{\def\LTcaptype{none} % do not increment counter
\begin{longtable}[]{@{}
  >{\raggedright\arraybackslash}p{(\linewidth - 24\tabcolsep) * \real{0.0746}}
  >{\raggedright\arraybackslash}p{(\linewidth - 24\tabcolsep) * \real{0.0746}}
  >{\raggedright\arraybackslash}p{(\linewidth - 24\tabcolsep) * \real{0.0746}}
  >{\raggedright\arraybackslash}p{(\linewidth - 24\tabcolsep) * \real{0.0746}}
  >{\raggedright\arraybackslash}p{(\linewidth - 24\tabcolsep) * \real{0.0746}}
  >{\raggedright\arraybackslash}p{(\linewidth - 24\tabcolsep) * \real{0.0746}}
  >{\raggedright\arraybackslash}p{(\linewidth - 24\tabcolsep) * \real{0.0746}}
  >{\raggedright\arraybackslash}p{(\linewidth - 24\tabcolsep) * \real{0.0746}}
  >{\raggedright\arraybackslash}p{(\linewidth - 24\tabcolsep) * \real{0.0746}}
  >{\raggedright\arraybackslash}p{(\linewidth - 24\tabcolsep) * \real{0.0746}}
  >{\raggedright\arraybackslash}p{(\linewidth - 24\tabcolsep) * \real{0.0746}}
  >{\raggedright\arraybackslash}p{(\linewidth - 24\tabcolsep) * \real{0.0746}}
  >{\raggedright\arraybackslash}p{(\linewidth - 24\tabcolsep) * \real{0.1045}}@{}}
\toprule\noalign{}
\begin{minipage}[b]{\linewidth}\raggedright
Jan
\end{minipage} & \begin{minipage}[b]{\linewidth}\raggedright
Fev
\end{minipage} & \begin{minipage}[b]{\linewidth}\raggedright
Mar
\end{minipage} & \begin{minipage}[b]{\linewidth}\raggedright
Abr
\end{minipage} & \begin{minipage}[b]{\linewidth}\raggedright
Mai
\end{minipage} & \begin{minipage}[b]{\linewidth}\raggedright
Jun
\end{minipage} & \begin{minipage}[b]{\linewidth}\raggedright
Jul
\end{minipage} & \begin{minipage}[b]{\linewidth}\raggedright
Ago
\end{minipage} & \begin{minipage}[b]{\linewidth}\raggedright
Set
\end{minipage} & \begin{minipage}[b]{\linewidth}\raggedright
Out
\end{minipage} & \begin{minipage}[b]{\linewidth}\raggedright
Nov
\end{minipage} & \begin{minipage}[b]{\linewidth}\raggedright
Dez
\end{minipage} & \begin{minipage}[b]{\linewidth}\raggedright
Média
\end{minipage} \\
\midrule\noalign{}
\endhead
\bottomrule\noalign{}
\endlastfoot
6.56 & 6.11 & 5.31 & 4.26 & 3.15 & 2.69 & 2.99 & 3.69 & 4.36 & 5.27 &
6.39 & 6.73 & \textbf{4.79} \\
\end{longtable}
}

\textbf{Inclinação solar recomendada:} 27° N (anual) · 37° N (inverno
jun-ago) · 17° N (verão nov-jan)

\paragraph{Precipitação (mm/mês)}\label{precipitauxe7uxe3o-mmmuxeas}

{\def\LTcaptype{none} % do not increment counter
\begin{longtable}[]{@{}
  >{\raggedright\arraybackslash}p{(\linewidth - 24\tabcolsep) * \real{0.0704}}
  >{\raggedright\arraybackslash}p{(\linewidth - 24\tabcolsep) * \real{0.0704}}
  >{\raggedright\arraybackslash}p{(\linewidth - 24\tabcolsep) * \real{0.0704}}
  >{\raggedright\arraybackslash}p{(\linewidth - 24\tabcolsep) * \real{0.0704}}
  >{\raggedright\arraybackslash}p{(\linewidth - 24\tabcolsep) * \real{0.0704}}
  >{\raggedright\arraybackslash}p{(\linewidth - 24\tabcolsep) * \real{0.0704}}
  >{\raggedright\arraybackslash}p{(\linewidth - 24\tabcolsep) * \real{0.0704}}
  >{\raggedright\arraybackslash}p{(\linewidth - 24\tabcolsep) * \real{0.0704}}
  >{\raggedright\arraybackslash}p{(\linewidth - 24\tabcolsep) * \real{0.0704}}
  >{\raggedright\arraybackslash}p{(\linewidth - 24\tabcolsep) * \real{0.0704}}
  >{\raggedright\arraybackslash}p{(\linewidth - 24\tabcolsep) * \real{0.0704}}
  >{\raggedright\arraybackslash}p{(\linewidth - 24\tabcolsep) * \real{0.0704}}
  >{\raggedright\arraybackslash}p{(\linewidth - 24\tabcolsep) * \real{0.1549}}@{}}
\toprule\noalign{}
\begin{minipage}[b]{\linewidth}\raggedright
Jan
\end{minipage} & \begin{minipage}[b]{\linewidth}\raggedright
Fev
\end{minipage} & \begin{minipage}[b]{\linewidth}\raggedright
Mar
\end{minipage} & \begin{minipage}[b]{\linewidth}\raggedright
Abr
\end{minipage} & \begin{minipage}[b]{\linewidth}\raggedright
Mai
\end{minipage} & \begin{minipage}[b]{\linewidth}\raggedright
Jun
\end{minipage} & \begin{minipage}[b]{\linewidth}\raggedright
Jul
\end{minipage} & \begin{minipage}[b]{\linewidth}\raggedright
Ago
\end{minipage} & \begin{minipage}[b]{\linewidth}\raggedright
Set
\end{minipage} & \begin{minipage}[b]{\linewidth}\raggedright
Out
\end{minipage} & \begin{minipage}[b]{\linewidth}\raggedright
Nov
\end{minipage} & \begin{minipage}[b]{\linewidth}\raggedright
Dez
\end{minipage} & \begin{minipage}[b]{\linewidth}\raggedright
Total/ano
\end{minipage} \\
\midrule\noalign{}
\endhead
\bottomrule\noalign{}
\endlastfoot
156 & 125 & 149 & 186 & 147 & 112 & 92 & 90 & 122 & 215 & 190 & 188 &
\textbf{1770 mm} \\
\end{longtable}
}

\paragraph{Poluição Luminosa}\label{poluiuxe7uxe3o-luminosa}

{\def\LTcaptype{none} % do not increment counter
\begin{longtable}[]{@{}ll@{}}
\toprule\noalign{}
Parâmetro & Valor \\
\midrule\noalign{}
\endhead
\bottomrule\noalign{}
\endlastfoot
Escala Bortle & 5 --- Céu suburbano \\
Radiância artificial & 5.0 nW/cm²/sr \\
\end{longtable}
}
\paragraph{Solo (SoilGrids 2.0, média ponderada 0--30
cm)}

{\def\LTcaptype{none} % do not increment counter
\begin{longtable}[]{@{}ll@{}}
\toprule\noalign{}
Parâmetro & Valor \\
\midrule\noalign{}
\endhead
\bottomrule\noalign{}
\endlastfoot
pH (H₂O) & 5.02 \\
Carbono orgânico (SOC) & 21.27 g/kg \\
Argila & 47.42 \% \\
Areia & 19.62 \% \\
Silte (calc.) & 33.0 \% \\
Densidade aparente & 1.24 g/cm³ \\
\textbf{Aptidão agrícola} & \textbf{Média} \\
\end{longtable}
}

Fonte: ISRIC SoilGrids 2.0 via WCS. Coords: -27.3333°, -55.8333°. Média
ponderada camadas 0-5, 5-15, 15-30 cm.

\begin{itemize}
\tightlist
\item
  \textbf{Agua:} Cobertura de rede de água e poços comunitários.
\item
  \textbf{Energia:} Rede ANDE estável.
\item
  \textbf{Producao:} Olericultura, pequenas granjas e serviços urbanos.
\end{itemize}
\subsubsection{Indicadores Sociais}

\textbf{Fontes:} PNUD 2020, INE Censo 2022, INE EPHC 2023, Ministerio
Público 2024\\
\textbf{Nota:} Valores marcados como \emph{(dept.)} referem-se ao
departamento de Itapúa; valores \emph{(dist.)} são específicos deste
distrito. Dados marcados \emph{(est.)} são estimativas.

{\def\LTcaptype{none} % do not increment counter
\begin{longtable}[]{@{}
  >{\raggedright\arraybackslash}p{(\linewidth - 4\tabcolsep) * \real{0.4231}}
  >{\raggedright\arraybackslash}p{(\linewidth - 4\tabcolsep) * \real{0.2692}}
  >{\raggedright\arraybackslash}p{(\linewidth - 4\tabcolsep) * \real{0.3077}}@{}}
\toprule\noalign{}
\begin{minipage}[b]{\linewidth}\raggedright
Indicador
\end{minipage} & \begin{minipage}[b]{\linewidth}\raggedright
Valor
\end{minipage} & \begin{minipage}[b]{\linewidth}\raggedright
Âmbito
\end{minipage} \\
\midrule\noalign{}
\endhead
\bottomrule\noalign{}
\endlastfoot
IDH (2020) & 0,728 & dept. \\
Ranking IDH nacional & 4/18 & dept. \\
Esperança de vida & 74,0 anos & dept. \\
Escolaridade média & 8,9 anos & dept. \\
RNB per capita (USD PPA) & 10.100 & dept. \\
Pobreza monetária (\%) & 22,1\% & dept. \\
Pobreza extrema (\%) & 4,8\% & dept. \\
Índice de Gini & 0,432 & dept. \\
Acesso a água potável (\%) & 84,6\% & dept. \\
Acesso a saneamento (\%) & 83,1\% & dept. \\
Taxa de homicídios (est., /100k hab) & \textasciitilde4--6 (estimativa,
próxima à média nacional) & dept. \\
Índice de segurança & médio & dept. \\
População (Censo 2022) & 47.717 & dist. \\
Idade mediana & N/D & dist. \\
Taxa de fecundidade & N/D & dist. \\
\end{longtable}
}
\subsubsection{Combustível}

\textbf{Referência:} PETROPAR / postos locais (2024) \textbf{Tipo de
localidade:} Interior

{\def\LTcaptype{none} % do not increment counter
\begin{longtable}[]{@{}lll@{}}
\toprule\noalign{}
Combustível & USD/litro & Gs/litro (aprox.) \\
\midrule\noalign{}
\endhead
\bottomrule\noalign{}
\endlastfoot
Gasolina 93 oct & 0.98 & 7,252 \\
Gasolina 97 oct (premium) & 1.08 & 7,992 \\
Diesel & 0.91 & 6,734 \\
\end{longtable}
}

\begin{quote}
Preços podem variar ±5\% conforme posto e sazonalidade. Chaco e interior
remoto apresentam maior variação.
\end{quote}

\subsubsection{Cobertura Celular}

\textbf{Fonte:} CONATEL PY / operadoras (2024)

{\def\LTcaptype{none} % do not increment counter
\begin{longtable}[]{@{}ll@{}}
\toprule\noalign{}
Parâmetro & Valor \\
\midrule\noalign{}
\endhead
\bottomrule\noalign{}
\endlastfoot
Cobertura 4G população (dept.) & 92\% \\
Cobertura 4G área rural & 82\% \\
Melhor operadora & Tigo \\
Qualidade rural & boa \\
\end{longtable}
}

\begin{quote}
Para áreas rurais fora do núcleo urbano, recomenda-se chip Tigo como
principal e Personal como backup.
\end{quote}

\subsubsection{Internet}

\textbf{Fonte:} CONATEL / Speedtest Ookla (2024)

{\def\LTcaptype{none} % do not increment counter
\begin{longtable}[]{@{}ll@{}}
\toprule\noalign{}
Parâmetro & Valor \\
\midrule\noalign{}
\endhead
\bottomrule\noalign{}
\endlastfoot
Velocidade média download & 75 Mbps \\
Domicílios com internet (dept.) & 78\% \\
Tecnologia predominante & fibra/rádio \\
Opção rural & Starlink disponível (\textasciitilde USD 44/mês) \\
\end{longtable}
}

\subsubsection{Mercado Imobiliário e Terra
Rural}

\textbf{Fonte:} INDERT / Clasificados.com.py (2024)

{\def\LTcaptype{none} % do not increment counter
\begin{longtable}[]{@{}ll@{}}
\toprule\noalign{}
Tipo & Referência \\
\midrule\noalign{}
\endhead
\bottomrule\noalign{}
\endlastfoot
Terra agrícola alta prod. (USD/ha) & 6,900 \\
Imóvel urbano (USD/m²) & 1,300 \\
Aluguel 2 quartos (USD/mês) & 550 \\
\end{longtable}
}

\begin{quote}
Valores de referência departamental. Localidades menores podem ter
preços 20--40\% abaixo da capital departamental.
\end{quote}

\subsubsection{Saúde}

\textbf{Fonte:} MSPBS / IPS Paraguay (2024-2026), consolidação
departamental e proxy local

{\def\LTcaptype{none} % do not increment counter
\begin{longtable}[]{@{}ll@{}}
\toprule\noalign{}
Serviço & Disponibilidade \\
\midrule\noalign{}
\endhead
\bottomrule\noalign{}
\endlastfoot
USF / Posto de Saúde & sim \\
Hospital Regional & sim \\
IPS (seguro social) & sim \\
Farmácia & sim \\
Distância ao hospital de referência & \textasciitilde18 km
(Encarnacion) \\
\end{longtable}
}

\textbf{Principais estabelecimentos:} USF local; referencia hospitalar
em Encarnacion

\textbf{Observação para imigrantes:} Atendimento primario local ou em
raio curto; casos de maior complexidade seguem para Encarnacion.
Cobertura privada continua recomendada para especialidades e urgências
de maior complexidade.
\newpage
\secaoDiagnostico{Capitan Meza}{\mapaDistrito{0704}}{Capitán Meza é um distrito agrícola promissor com solos de excelente
qualidade. Sua localização estratégica no Corredor da Exportação oferece
vantagens logísticas futuras, mas sua baixa densidade e foco em
commodities exigem planejamento de autossuficiência.

\subsubsection{Por que sim}

\begin{itemize}
\tightlist
\item
  Solo de altíssima produtividade (Terra Roxa).
\item
  Conectividade em franca expansão (Corredor da Exportação).
\end{itemize}

\subsubsection{Por que não}

\begin{itemize}
\tightlist
\item
  Dependência excessiva de infraestrutura externa para serviços
  especializados.
\item
  Vulnerabilidade a fluxos de fronteira.
\end{itemize}

\subsubsection{Recomendação
Técnica}

Recomendado para produtores agrícolas mecanizados e para quem busca uma
área rural produtiva e segura no eixo do Rio Paraná.

\subsubsection{Analise de Sensibilidade}

\begin{itemize}
\tightlist
\item
  Cenario otimista: GSS 6.3 (Instalação de indústrias de processamento
  local).
\item
  Cenario conservador: GSS 5.3 (Crise hídrica afetando o transporte
  fluvial no Paraná).
\end{itemize}

\subsubsection{}

\begin{itemize}
\tightlist
\item
  Censo 2022 INE.
\item
  Planos de expansão rodoviária MOPC.
\item
  Relatórios de exportação portuária.
\end{itemize}}{dist:07_Itapua:Capitan_Meza}
\begin{table}[ht!]
\small \begin{tabular}{p{3.0cm}p{0.8cm}p{6.0cm}} \toprule
\textbf{Critério} & \textbf{Nota} & \textbf{Justificativa} \\ \midrule
A - Ameaças Estratégicas & 4.6 & - \\
B - Risco Social & 5.5 & - \\
C - Riscos Naturais & 7.0 & - \\
D - Autossuficiência & 6.5 & - \\
E - Institucional & 6.0 & - \\
\midrule \textbf{GSS FINAL} & \textbf{5.8} & \textbf{Seguro (Potencial moderado de resiliência).} \\ \bottomrule \end{tabular}
\end{table}
\subsubsection{Vulnerabilidades e
Mitigação}\label{vuln-Capitan_Meza-vulnerabilidades-e-mitigauxe7uxe3o}

\begin{itemize}
\item
  \textbf{Dependência Logística:} Grande foco em exportação pode gerar
  crises locais em bloqueios de portos. Mitigação: Reservas financeiras
  e diversificação para mercado interno.
\item
  \textbf{Fronteira:} Exposição a atividades de contrabando na área
  ribeirinha. Mitigação: Fortalecimento da vigilância patrimonial.
\item
  INE Censo 2022: 9.253 habitantes.
\item
  Preço da Terra: US\$ 4.000 a US\$ 6.000 por hectare.
\item
  Homicídios: Baixa incidência regional.
\end{itemize}
\subsubsection{Dossiê de Campo}\label{dossie-Capitan_Meza-dossiuxea-de-campo}
\begin{description}[style=nextline,leftmargin=0.5cm,font=\bfseries]
\item[Coordenadas]  26°40′S 55°14′W.
\item[Topografia]  Relevo ondulado com descidas acentuadas para o Rio Paraná. Área de \textasciitilde 440 km².
\item[Alvos Estrategicos]  Centro de produção agrícola e porto graneleiro secundário. Proximidade com a fronteira argentina. Sem alvos militares.
\item[Fallout]  Ventos predominantes do Norte.
\item[Populacao]  9.253 (Censo 2022 Final).
\item[Densidade]  \textasciitilde 21,0 hab/km².
\item[Vias de Acesso]  Acesso pela Ruta PY06 (via conexão asfáltica) e pela "Ruta del Río" (Projeto Corredor da Exportação).
\item[Servicos]  Infraestrutura de saúde e educação básicas. Dependência de centros como Obligado ou Encarnación para serviços complexos.
\item[Hidrologia]  Diversos arroios. Risco de alagamentos pontuais em áreas ribeirinhas do Rio Paraná.
\item[Sismicidade]  Baixa.
\item[Solo]  Terra Roxa de origem basáltica, altamente fértil.
\item[Agua]  Poços artesianos e sistemas locais de distribuição.
\item[Energia]  Rede ANDE.
\item[Producao]  Soja, milho, trigo, girassol e erva-mate.
\item[Seguranca]  Região de baixa criminalidade, característica de zona agrícola produtiva.
\item[Clima Social]  Estabilidade social elevada, presença de descendentes de imigrantes (cultura de trabalho).
\end{description}
\textbf{Fonte climática:} NASA POWER Climatology API (período 2001-2020)

\textbf{Fonte luminosa:} estimativa world atlas

\paragraph{Irradiação Solar
(kWh/m²/dia)}\label{irradiauxe7uxe3o-solar-kwhmuxb2dia}

{\def\LTcaptype{none} % do not increment counter
\begin{longtable}[]{@{}
  >{\raggedright\arraybackslash}p{(\linewidth - 24\tabcolsep) * \real{0.0746}}
  >{\raggedright\arraybackslash}p{(\linewidth - 24\tabcolsep) * \real{0.0746}}
  >{\raggedright\arraybackslash}p{(\linewidth - 24\tabcolsep) * \real{0.0746}}
  >{\raggedright\arraybackslash}p{(\linewidth - 24\tabcolsep) * \real{0.0746}}
  >{\raggedright\arraybackslash}p{(\linewidth - 24\tabcolsep) * \real{0.0746}}
  >{\raggedright\arraybackslash}p{(\linewidth - 24\tabcolsep) * \real{0.0746}}
  >{\raggedright\arraybackslash}p{(\linewidth - 24\tabcolsep) * \real{0.0746}}
  >{\raggedright\arraybackslash}p{(\linewidth - 24\tabcolsep) * \real{0.0746}}
  >{\raggedright\arraybackslash}p{(\linewidth - 24\tabcolsep) * \real{0.0746}}
  >{\raggedright\arraybackslash}p{(\linewidth - 24\tabcolsep) * \real{0.0746}}
  >{\raggedright\arraybackslash}p{(\linewidth - 24\tabcolsep) * \real{0.0746}}
  >{\raggedright\arraybackslash}p{(\linewidth - 24\tabcolsep) * \real{0.0746}}
  >{\raggedright\arraybackslash}p{(\linewidth - 24\tabcolsep) * \real{0.1045}}@{}}
\toprule\noalign{}
\begin{minipage}[b]{\linewidth}\raggedright
Jan
\end{minipage} & \begin{minipage}[b]{\linewidth}\raggedright
Fev
\end{minipage} & \begin{minipage}[b]{\linewidth}\raggedright
Mar
\end{minipage} & \begin{minipage}[b]{\linewidth}\raggedright
Abr
\end{minipage} & \begin{minipage}[b]{\linewidth}\raggedright
Mai
\end{minipage} & \begin{minipage}[b]{\linewidth}\raggedright
Jun
\end{minipage} & \begin{minipage}[b]{\linewidth}\raggedright
Jul
\end{minipage} & \begin{minipage}[b]{\linewidth}\raggedright
Ago
\end{minipage} & \begin{minipage}[b]{\linewidth}\raggedright
Set
\end{minipage} & \begin{minipage}[b]{\linewidth}\raggedright
Out
\end{minipage} & \begin{minipage}[b]{\linewidth}\raggedright
Nov
\end{minipage} & \begin{minipage}[b]{\linewidth}\raggedright
Dez
\end{minipage} & \begin{minipage}[b]{\linewidth}\raggedright
Média
\end{minipage} \\
\midrule\noalign{}
\endhead
\bottomrule\noalign{}
\endlastfoot
6.56 & 6.06 & 5.32 & 4.33 & 3.21 & 2.76 & 3.08 & 3.81 & 4.39 & 5.25 &
6.31 & 6.67 & \textbf{4.81} \\
\end{longtable}
}

\textbf{Inclinação solar recomendada:} 27° N (anual) · 37° N (inverno
jun-ago) · 17° N (verão nov-jan)

\paragraph{Precipitação (mm/mês)}\label{precipitauxe7uxe3o-mmmuxeas}

{\def\LTcaptype{none} % do not increment counter
\begin{longtable}[]{@{}
  >{\raggedright\arraybackslash}p{(\linewidth - 24\tabcolsep) * \real{0.0704}}
  >{\raggedright\arraybackslash}p{(\linewidth - 24\tabcolsep) * \real{0.0704}}
  >{\raggedright\arraybackslash}p{(\linewidth - 24\tabcolsep) * \real{0.0704}}
  >{\raggedright\arraybackslash}p{(\linewidth - 24\tabcolsep) * \real{0.0704}}
  >{\raggedright\arraybackslash}p{(\linewidth - 24\tabcolsep) * \real{0.0704}}
  >{\raggedright\arraybackslash}p{(\linewidth - 24\tabcolsep) * \real{0.0704}}
  >{\raggedright\arraybackslash}p{(\linewidth - 24\tabcolsep) * \real{0.0704}}
  >{\raggedright\arraybackslash}p{(\linewidth - 24\tabcolsep) * \real{0.0704}}
  >{\raggedright\arraybackslash}p{(\linewidth - 24\tabcolsep) * \real{0.0704}}
  >{\raggedright\arraybackslash}p{(\linewidth - 24\tabcolsep) * \real{0.0704}}
  >{\raggedright\arraybackslash}p{(\linewidth - 24\tabcolsep) * \real{0.0704}}
  >{\raggedright\arraybackslash}p{(\linewidth - 24\tabcolsep) * \real{0.0704}}
  >{\raggedright\arraybackslash}p{(\linewidth - 24\tabcolsep) * \real{0.1549}}@{}}
\toprule\noalign{}
\begin{minipage}[b]{\linewidth}\raggedright
Jan
\end{minipage} & \begin{minipage}[b]{\linewidth}\raggedright
Fev
\end{minipage} & \begin{minipage}[b]{\linewidth}\raggedright
Mar
\end{minipage} & \begin{minipage}[b]{\linewidth}\raggedright
Abr
\end{minipage} & \begin{minipage}[b]{\linewidth}\raggedright
Mai
\end{minipage} & \begin{minipage}[b]{\linewidth}\raggedright
Jun
\end{minipage} & \begin{minipage}[b]{\linewidth}\raggedright
Jul
\end{minipage} & \begin{minipage}[b]{\linewidth}\raggedright
Ago
\end{minipage} & \begin{minipage}[b]{\linewidth}\raggedright
Set
\end{minipage} & \begin{minipage}[b]{\linewidth}\raggedright
Out
\end{minipage} & \begin{minipage}[b]{\linewidth}\raggedright
Nov
\end{minipage} & \begin{minipage}[b]{\linewidth}\raggedright
Dez
\end{minipage} & \begin{minipage}[b]{\linewidth}\raggedright
Total/ano
\end{minipage} \\
\midrule\noalign{}
\endhead
\bottomrule\noalign{}
\endlastfoot
153 & 129 & 137 & 185 & 155 & 120 & 99 & 84 & 117 & 206 & 185 & 182 &
\textbf{1751 mm} \\
\end{longtable}
}

\paragraph{Poluição Luminosa}\label{poluiuxe7uxe3o-luminosa}

{\def\LTcaptype{none} % do not increment counter
\begin{longtable}[]{@{}ll@{}}
\toprule\noalign{}
Parâmetro & Valor \\
\midrule\noalign{}
\endhead
\bottomrule\noalign{}
\endlastfoot
Escala Bortle & 5 --- Céu suburbano \\
Radiância artificial & 5.0 nW/cm²/sr \\
\end{longtable}
}
\paragraph{Solo (SoilGrids 2.0, média ponderada 0--30
cm)}

{\def\LTcaptype{none} % do not increment counter
\begin{longtable}[]{@{}ll@{}}
\toprule\noalign{}
Parâmetro & Valor \\
\midrule\noalign{}
\endhead
\bottomrule\noalign{}
\endlastfoot
pH (H₂O) & 5.28 \\
Carbono orgânico (SOC) & 19.21 g/kg \\
Argila & 40.07 \% \\
Areia & 29.43 \% \\
Silte (calc.) & 30.5 \% \\
Densidade aparente & 1.27 g/cm³ \\
\textbf{Aptidão agrícola} & \textbf{Média} \\
\end{longtable}
}

Fonte: ISRIC SoilGrids 2.0 via WCS. Coords: -26.6667°, -55.2333°. Média
ponderada camadas 0-5, 5-15, 15-30 cm.

\begin{itemize}
\tightlist
\item
  \textbf{Agua:} Poços artesianos e sistemas locais de distribuição.
\item
  \textbf{Energia:} Rede ANDE.
\item
  \textbf{Producao:} Soja, milho, trigo, girassol e erva-mate.
\end{itemize}
\subsubsection{Indicadores Sociais}

\textbf{Fontes:} PNUD 2020, INE Censo 2022, INE EPHC 2023, Ministerio
Público 2024\\
\textbf{Nota:} Valores marcados como \emph{(dept.)} referem-se ao
departamento de Itapúa; valores \emph{(dist.)} são específicos deste
distrito. Dados marcados \emph{(est.)} são estimativas.

{\def\LTcaptype{none} % do not increment counter
\begin{longtable}[]{@{}
  >{\raggedright\arraybackslash}p{(\linewidth - 4\tabcolsep) * \real{0.4231}}
  >{\raggedright\arraybackslash}p{(\linewidth - 4\tabcolsep) * \real{0.2692}}
  >{\raggedright\arraybackslash}p{(\linewidth - 4\tabcolsep) * \real{0.3077}}@{}}
\toprule\noalign{}
\begin{minipage}[b]{\linewidth}\raggedright
Indicador
\end{minipage} & \begin{minipage}[b]{\linewidth}\raggedright
Valor
\end{minipage} & \begin{minipage}[b]{\linewidth}\raggedright
Âmbito
\end{minipage} \\
\midrule\noalign{}
\endhead
\bottomrule\noalign{}
\endlastfoot
IDH (2020) & 0,728 & dept. \\
Ranking IDH nacional & 4/18 & dept. \\
Esperança de vida & 74,0 anos & dept. \\
Escolaridade média & 7.8 anos & dist. \\
RNB per capita (USD PPA) & 10.100 & dept. \\
Pobreza monetária (\%) & 22,1\% & dept. \\
Pobreza extrema (\%) & 4,8\% & dept. \\
Índice de Gini & 0,432 & dept. \\
Acesso a água potável (\%) & 84,6\% & dept. \\
Acesso a saneamento (\%) & 83,1\% & dept. \\
Taxa de homicídios (est., /100k hab) & \textasciitilde4--6 (estimativa,
próxima à média nacional) & dept. \\
Índice de segurança & médio & dept. \\
População (Censo 2022) & 9.253 hab. & dist. \\
Idade mediana & 29.0 anos & dist. \\
Taxa de fecundidade & N/D & dist. \\
\end{longtable}
}
\subsubsection{Combustível}

\textbf{Referência:} PETROPAR / postos locais (2024) \textbf{Tipo de
localidade:} Interior

{\def\LTcaptype{none} % do not increment counter
\begin{longtable}[]{@{}lll@{}}
\toprule\noalign{}
Combustível & USD/litro & Gs/litro (aprox.) \\
\midrule\noalign{}
\endhead
\bottomrule\noalign{}
\endlastfoot
Gasolina 93 oct & 0.98 & 7,252 \\
Gasolina 97 oct (premium) & 1.08 & 7,992 \\
Diesel & 0.91 & 6,734 \\
\end{longtable}
}

\begin{quote}
Preços podem variar ±5\% conforme posto e sazonalidade. Chaco e interior
remoto apresentam maior variação.
\end{quote}

\subsubsection{Cobertura Celular}

\textbf{Fonte:} CONATEL PY / operadoras (2024)

{\def\LTcaptype{none} % do not increment counter
\begin{longtable}[]{@{}ll@{}}
\toprule\noalign{}
Parâmetro & Valor \\
\midrule\noalign{}
\endhead
\bottomrule\noalign{}
\endlastfoot
Cobertura 4G população (dept.) & 92\% \\
Cobertura 4G área rural & 82\% \\
Melhor operadora & Tigo \\
Qualidade rural & boa \\
\end{longtable}
}

\begin{quote}
Para áreas rurais fora do núcleo urbano, recomenda-se chip Tigo como
principal e Personal como backup.
\end{quote}

\subsubsection{Internet}

\textbf{Fonte:} CONATEL / Speedtest Ookla (2024)

{\def\LTcaptype{none} % do not increment counter
\begin{longtable}[]{@{}ll@{}}
\toprule\noalign{}
Parâmetro & Valor \\
\midrule\noalign{}
\endhead
\bottomrule\noalign{}
\endlastfoot
Velocidade média download & 75 Mbps \\
Domicílios com internet (dept.) & 78\% \\
Tecnologia predominante & fibra/rádio \\
Opção rural & Starlink disponível (\textasciitilde USD 44/mês) \\
\end{longtable}
}

\subsubsection{Mercado Imobiliário e Terra
Rural}

\textbf{Fonte:} INDERT / Clasificados.com.py (2024)

{\def\LTcaptype{none} % do not increment counter
\begin{longtable}[]{@{}ll@{}}
\toprule\noalign{}
Tipo & Referência \\
\midrule\noalign{}
\endhead
\bottomrule\noalign{}
\endlastfoot
Terra agrícola alta prod. (USD/ha) & 6,900 \\
Imóvel urbano (USD/m²) & 1,300 \\
Aluguel 2 quartos (USD/mês) & 550 \\
\end{longtable}
}

\begin{quote}
Valores de referência departamental. Localidades menores podem ter
preços 20--40\% abaixo da capital departamental.
\end{quote}

\subsubsection{Saúde}

\textbf{Fonte:} MSPBS / IPS Paraguay (2024-2026), consolidação
departamental e proxy local

{\def\LTcaptype{none} % do not increment counter
\begin{longtable}[]{@{}ll@{}}
\toprule\noalign{}
Serviço & Disponibilidade \\
\midrule\noalign{}
\endhead
\bottomrule\noalign{}
\endlastfoot
USF / Posto de Saúde & sim \\
Hospital Regional & não \\
IPS (seguro social) & sim \\
Farmácia & sim \\
Distância ao hospital de referência & \textasciitilde96 km
(Encarnacion) \\
\end{longtable}
}

\textbf{Principais estabelecimentos:} USF local; referencia hospitalar
em Encarnacion

\textbf{Observação para imigrantes:} Atendimento primario local ou em
raio curto; casos de maior complexidade seguem para Encarnacion.
Cobertura privada continua recomendada para especialidades e urgências
de maior complexidade.
\newpage
\secaoDiagnostico{Capitan Miranda}{\mapaDistrito{0705}}{Capitan Miranda (Itapua) apresenta perfil operacional consolidado para
comparacao territorial, com leitura conjunta de infraestrutura, risco
hidroclimatico e capacidade de continuidade de servicos essenciais.

\subsubsection{Por que sim}

\begin{itemize}
\tightlist
\item
  Base institucional de referencia disponivel e rastreavel.
\item
  Estrutura comparativa (A-E/GSS) consistente para decisao relativa.
\item
  Plano de mitigacao acionavel ja definido no dossie.
\end{itemize}

\subsubsection{Por que não}

\begin{itemize}
\tightlist
\item
  Serie oficial distrital granular ainda pode apresentar lacunas
  temporais.
\item
  Sensibilidade do score exige monitoramento continuo de risco e
  conectividade.
\item
  Decisao final depende de validacao de campo e janela temporal recente.
\end{itemize}

\subsubsection{Pre-condicoes Minimas de
Decisao}

\begin{itemize}
\tightlist
\item
  Atualizar indicadores criticos em ciclo trimestral.
\item
  Confirmar trafegabilidade e contingencia hidrica/energetica local.
\item
  Validar checklist de resiliencia domiciliar/logistica antes de decisao
  final.
\end{itemize}

\subsubsection{Recomendação
Técnica}

Uso recomendado como candidato em analise multicriterio, condicionado ao
cumprimento das pre-condicoes acima. Referencia atual de score: GSS 5.2;
risco predominante: alto.

\subsubsection{}

\begin{itemize}
\tightlist
\item
  \begin{enumerate}
  \def\labelenumi{\arabic{enumi})}
  \tightlist
  \item
    Delimitacao territorial validada por georreferencia institucional
    (Fonte: acesso: 2026-03-05).
  \end{enumerate}
\item
  \begin{enumerate}
  \def\labelenumi{\arabic{enumi})}
  \setcounter{enumi}{1}
  \tightlist
  \item
    Populacao municipal/distrital referenciada por base censitaria
    nacional (Fonte: acesso: 2026-03-05).
  \end{enumerate}
\item
  \begin{enumerate}
  \def\labelenumi{\arabic{enumi})}
  \setcounter{enumi}{2}
  \tightlist
  \item
    Comparabilidade intra-departamental apoiada em indicadores
    distritais oficiais (Fonte: acesso: 2026-03-05).
  \end{enumerate}
\item
  \begin{enumerate}
  \def\labelenumi{\arabic{enumi})}
  \setcounter{enumi}{3}
  \tightlist
  \item
    Estado macro de conectividade viaria ancorado em informacoes de
    rodovias (Fonte: acesso: 2026-03-05).
  \end{enumerate}
\item
  \begin{enumerate}
  \def\labelenumi{\arabic{enumi})}
  \setcounter{enumi}{4}
  \tightlist
  \item
    Exposicao hidrometeorologica monitorada por avisos oficiais (Fonte:
    acesso: 2026-03-05).
  \end{enumerate}
\item
  \begin{enumerate}
  \def\labelenumi{\arabic{enumi})}
  \setcounter{enumi}{5}
  \tightlist
  \item
    Historico de contexto meteorologico apoiado em publicacoes tecnicas
    (Fonte: acesso: 2026-03-05).
  \end{enumerate}
\item
  \begin{enumerate}
  \def\labelenumi{\arabic{enumi})}
  \setcounter{enumi}{6}
  \tightlist
  \item
    Capacidade institucional de resposta a eventos apoiada em acoes da
    SEN (Fonte: acesso: 2026-03-05).
  \end{enumerate}
\item
  \begin{enumerate}
  \def\labelenumi{\arabic{enumi})}
  \setcounter{enumi}{7}
  \tightlist
  \item
    Infraestrutura energetica analisada via operador nacional (Fonte:
    acesso: 2026-03-05).
  \end{enumerate}
\item
  \begin{enumerate}
  \def\labelenumi{\arabic{enumi})}
  \setcounter{enumi}{8}
  \tightlist
  \item
    Referencias de obras e logistica nacional verificadas no MOPC
    (Fonte: acesso: 2026-03-05).
  \end{enumerate}
\item
  \begin{enumerate}
  \def\labelenumi{\arabic{enumi})}
  \setcounter{enumi}{9}
  \tightlist
  \item
    Contexto sociopolitico institucional referenciado por autoridade
    eleitoral (Fonte: acesso: 2026-03-05).
  \end{enumerate}
\item
  \begin{enumerate}
  \def\labelenumi{\arabic{enumi})}
  \setcounter{enumi}{10}
  \tightlist
  \item
    Camada cartografica de apoio eleitoral/territorial consultada
    (Fonte: acesso: 2026-03-05).
  \end{enumerate}
\item
  \begin{enumerate}
  \def\labelenumi{\arabic{enumi})}
  \setcounter{enumi}{11}
  \tightlist
  \item
    Dados publicos complementares para verificacao cruzada (Fonte:
    acesso: 2026-03-05).
  \end{enumerate}
\end{itemize}}{dist:07_Itapua:Capitan_Miranda}
\begin{table}[ht!]
\small \begin{tabular}{p{3.0cm}p{0.8cm}p{6.0cm}} \toprule
\textbf{Critério} & \textbf{Nota} & \textbf{Justificativa} \\ \midrule
A - Ameaças Estratégicas & 5.5 & - \\
B - Risco Social & 6.2 & - \\
C - Riscos Naturais & 3.1 & - \\
D - Autossuficiência & 4.4 & - \\
E - Institucional & 6.1 & - \\
\midrule \textbf{GSS FINAL} & \textbf{5.2} & \textbf{Sensivel (requer mitigacao robusta).} \\ \bottomrule \end{tabular}
\end{table}
\subsubsection{Vulnerabilidades e
Mitigação}\label{vuln-Capitan_Miranda-vulnerabilidades-e-mitigauxe7uxe3o}

\begin{itemize}
\tightlist
\item
  Alta exposicao hidrometeorologica (\textgreater=65/100).
\item
  Conectividade viaria baixa (\textless50/100).
\item
  Estoque de contingencia para 14 dias (agua/alimentos/combustivel),
  priorizado quando risco hidro=86/100.
\item
  Duplicar rotas/logistica quando conectividade viaria=37/100 for
  \textless70; manter rota secundaria mapeada e testada trimestralmente.
\item
  Plano de energia resiliente com autonomia minima de 72h
  (geracao/backup), revisao semestral.
\end{itemize}

Notas (0-10): - A: 5.5 - B: 6.2 - C: 3.1 - D: 4.4 - E: 6.1

Rastreabilidade das notas (base nos indicadores da secao 9): - A
(geoestrategia): combina conectividade viaria (37/100) e exposicao hidro
(86/100). - B (infraestrutura): combina conectividade (37/100) e escala
populacional (141082 hab). - C (riscos naturais): derivada da exposicao
hidrometeorologica (86/100). - D (autossuficiencia): combina
conectividade, ajuste de densidade (506.0 hab/km2) e risco hidro. - E
(ambiente sociopolitico): combinacao conservadora de conectividade,
escala populacional e risco hidro.

Calculo: - GSS = ((A x 2.5) + (B x 2.0) + (C x 1.5) + (D x 2.0) + (E x
2.0)) / 10 - GSS: 5.2

Classificacao: - Sensivel (requer mitigacao robusta).

Sintese aplicada: - Dados textuais consolidados com base nas fontes
oficiais acima; proximos ciclos podem aprofundar indicadores numericos
especificos por distrito.
\subsubsection{Dossiê de Campo}\label{dossie-Capitan_Miranda-dossiuxea-de-campo}
\begin{description}[style=nextline,leftmargin=0.5cm,font=\bfseries]
\item[Coordenadas]  27°13′S 55°48′W (geocodificado via Nominatim).
\end{description}
\textbf{Fonte climática:} NASA POWER Climatology API (período 2001-2020)

\textbf{Fonte luminosa:} estimativa world atlas

\paragraph{Irradiação Solar
(kWh/m²/dia)}\label{irradiauxe7uxe3o-solar-kwhmuxb2dia}

{\def\LTcaptype{none} % do not increment counter
\begin{longtable}[]{@{}
  >{\raggedright\arraybackslash}p{(\linewidth - 24\tabcolsep) * \real{0.0746}}
  >{\raggedright\arraybackslash}p{(\linewidth - 24\tabcolsep) * \real{0.0746}}
  >{\raggedright\arraybackslash}p{(\linewidth - 24\tabcolsep) * \real{0.0746}}
  >{\raggedright\arraybackslash}p{(\linewidth - 24\tabcolsep) * \real{0.0746}}
  >{\raggedright\arraybackslash}p{(\linewidth - 24\tabcolsep) * \real{0.0746}}
  >{\raggedright\arraybackslash}p{(\linewidth - 24\tabcolsep) * \real{0.0746}}
  >{\raggedright\arraybackslash}p{(\linewidth - 24\tabcolsep) * \real{0.0746}}
  >{\raggedright\arraybackslash}p{(\linewidth - 24\tabcolsep) * \real{0.0746}}
  >{\raggedright\arraybackslash}p{(\linewidth - 24\tabcolsep) * \real{0.0746}}
  >{\raggedright\arraybackslash}p{(\linewidth - 24\tabcolsep) * \real{0.0746}}
  >{\raggedright\arraybackslash}p{(\linewidth - 24\tabcolsep) * \real{0.0746}}
  >{\raggedright\arraybackslash}p{(\linewidth - 24\tabcolsep) * \real{0.0746}}
  >{\raggedright\arraybackslash}p{(\linewidth - 24\tabcolsep) * \real{0.1045}}@{}}
\toprule\noalign{}
\begin{minipage}[b]{\linewidth}\raggedright
Jan
\end{minipage} & \begin{minipage}[b]{\linewidth}\raggedright
Fev
\end{minipage} & \begin{minipage}[b]{\linewidth}\raggedright
Mar
\end{minipage} & \begin{minipage}[b]{\linewidth}\raggedright
Abr
\end{minipage} & \begin{minipage}[b]{\linewidth}\raggedright
Mai
\end{minipage} & \begin{minipage}[b]{\linewidth}\raggedright
Jun
\end{minipage} & \begin{minipage}[b]{\linewidth}\raggedright
Jul
\end{minipage} & \begin{minipage}[b]{\linewidth}\raggedright
Ago
\end{minipage} & \begin{minipage}[b]{\linewidth}\raggedright
Set
\end{minipage} & \begin{minipage}[b]{\linewidth}\raggedright
Out
\end{minipage} & \begin{minipage}[b]{\linewidth}\raggedright
Nov
\end{minipage} & \begin{minipage}[b]{\linewidth}\raggedright
Dez
\end{minipage} & \begin{minipage}[b]{\linewidth}\raggedright
Média
\end{minipage} \\
\midrule\noalign{}
\endhead
\bottomrule\noalign{}
\endlastfoot
6.56 & 6.11 & 5.31 & 4.26 & 3.15 & 2.69 & 2.99 & 3.69 & 4.36 & 5.27 &
6.39 & 6.73 & \textbf{4.79} \\
\end{longtable}
}

\textbf{Inclinação solar recomendada:} 27° N (anual) · 37° N (inverno
jun-ago) · 17° N (verão nov-jan)

\paragraph{Precipitação (mm/mês)}\label{precipitauxe7uxe3o-mmmuxeas}

{\def\LTcaptype{none} % do not increment counter
\begin{longtable}[]{@{}
  >{\raggedright\arraybackslash}p{(\linewidth - 24\tabcolsep) * \real{0.0704}}
  >{\raggedright\arraybackslash}p{(\linewidth - 24\tabcolsep) * \real{0.0704}}
  >{\raggedright\arraybackslash}p{(\linewidth - 24\tabcolsep) * \real{0.0704}}
  >{\raggedright\arraybackslash}p{(\linewidth - 24\tabcolsep) * \real{0.0704}}
  >{\raggedright\arraybackslash}p{(\linewidth - 24\tabcolsep) * \real{0.0704}}
  >{\raggedright\arraybackslash}p{(\linewidth - 24\tabcolsep) * \real{0.0704}}
  >{\raggedright\arraybackslash}p{(\linewidth - 24\tabcolsep) * \real{0.0704}}
  >{\raggedright\arraybackslash}p{(\linewidth - 24\tabcolsep) * \real{0.0704}}
  >{\raggedright\arraybackslash}p{(\linewidth - 24\tabcolsep) * \real{0.0704}}
  >{\raggedright\arraybackslash}p{(\linewidth - 24\tabcolsep) * \real{0.0704}}
  >{\raggedright\arraybackslash}p{(\linewidth - 24\tabcolsep) * \real{0.0704}}
  >{\raggedright\arraybackslash}p{(\linewidth - 24\tabcolsep) * \real{0.0704}}
  >{\raggedright\arraybackslash}p{(\linewidth - 24\tabcolsep) * \real{0.1549}}@{}}
\toprule\noalign{}
\begin{minipage}[b]{\linewidth}\raggedright
Jan
\end{minipage} & \begin{minipage}[b]{\linewidth}\raggedright
Fev
\end{minipage} & \begin{minipage}[b]{\linewidth}\raggedright
Mar
\end{minipage} & \begin{minipage}[b]{\linewidth}\raggedright
Abr
\end{minipage} & \begin{minipage}[b]{\linewidth}\raggedright
Mai
\end{minipage} & \begin{minipage}[b]{\linewidth}\raggedright
Jun
\end{minipage} & \begin{minipage}[b]{\linewidth}\raggedright
Jul
\end{minipage} & \begin{minipage}[b]{\linewidth}\raggedright
Ago
\end{minipage} & \begin{minipage}[b]{\linewidth}\raggedright
Set
\end{minipage} & \begin{minipage}[b]{\linewidth}\raggedright
Out
\end{minipage} & \begin{minipage}[b]{\linewidth}\raggedright
Nov
\end{minipage} & \begin{minipage}[b]{\linewidth}\raggedright
Dez
\end{minipage} & \begin{minipage}[b]{\linewidth}\raggedright
Total/ano
\end{minipage} \\
\midrule\noalign{}
\endhead
\bottomrule\noalign{}
\endlastfoot
154 & 122 & 146 & 184 & 150 & 110 & 90 & 82 & 113 & 206 & 193 & 180 &
\textbf{1732 mm} \\
\end{longtable}
}

\paragraph{Poluição Luminosa}\label{poluiuxe7uxe3o-luminosa}

{\def\LTcaptype{none} % do not increment counter
\begin{longtable}[]{@{}ll@{}}
\toprule\noalign{}
Parâmetro & Valor \\
\midrule\noalign{}
\endhead
\bottomrule\noalign{}
\endlastfoot
Escala Bortle & 5 --- Céu suburbano \\
Radiância artificial & 5.0 nW/cm²/sr \\
\end{longtable}
}
\paragraph{Solo (SoilGrids 2.0, média ponderada 0--30
cm)}

{\def\LTcaptype{none} % do not increment counter
\begin{longtable}[]{@{}ll@{}}
\toprule\noalign{}
Parâmetro & Valor \\
\midrule\noalign{}
\endhead
\bottomrule\noalign{}
\endlastfoot
pH (H₂O) & 5.17 \\
Carbono orgânico (SOC) & 20.93 g/kg \\
Argila & 45.11 \% \\
Areia & 17.6 \% \\
Silte (calc.) & 37.3 \% \\
Densidade aparente & 1.21 g/cm³ \\
\textbf{Aptidão agrícola} & \textbf{Média} \\
\end{longtable}
}

Fonte: ISRIC SoilGrids 2.0 via WCS. Coords: -27.2167°, -55.8°. Média
ponderada camadas 0-5, 5-15, 15-30 cm.
\subsubsection{Indicadores Sociais}

\textbf{Fontes:} PNUD 2020, INE Censo 2022, INE EPHC 2023, Ministerio
Público 2024\\
\textbf{Nota:} Valores marcados como \emph{(dept.)} referem-se ao
departamento de Itapúa; valores \emph{(dist.)} são específicos deste
distrito. Dados marcados \emph{(est.)} são estimativas.

{\def\LTcaptype{none} % do not increment counter
\begin{longtable}[]{@{}
  >{\raggedright\arraybackslash}p{(\linewidth - 4\tabcolsep) * \real{0.4231}}
  >{\raggedright\arraybackslash}p{(\linewidth - 4\tabcolsep) * \real{0.2692}}
  >{\raggedright\arraybackslash}p{(\linewidth - 4\tabcolsep) * \real{0.3077}}@{}}
\toprule\noalign{}
\begin{minipage}[b]{\linewidth}\raggedright
Indicador
\end{minipage} & \begin{minipage}[b]{\linewidth}\raggedright
Valor
\end{minipage} & \begin{minipage}[b]{\linewidth}\raggedright
Âmbito
\end{minipage} \\
\midrule\noalign{}
\endhead
\bottomrule\noalign{}
\endlastfoot
IDH (2020) & 0,728 & dept. \\
Ranking IDH nacional & 4/18 & dept. \\
Esperança de vida & 74,0 anos & dept. \\
Escolaridade média & 9.0 anos & dist. \\
RNB per capita (USD PPA) & 10.100 & dept. \\
Pobreza monetária (\%) & 22,1\% & dept. \\
Pobreza extrema (\%) & 4,8\% & dept. \\
Índice de Gini & 0,432 & dept. \\
Acesso a água potável (\%) & 84,6\% & dept. \\
Acesso a saneamento (\%) & 83,1\% & dept. \\
Taxa de homicídios (est., /100k hab) & \textasciitilde4--6 (estimativa,
próxima à média nacional) & dept. \\
Índice de segurança & médio & dept. \\
População (Censo 2022) & 9.528 hab. & dist. \\
Idade mediana & 31.0 anos & dist. \\
Taxa de fecundidade & N/D & dist. \\
\end{longtable}
}
\subsubsection{Combustível}

\textbf{Referência:} PETROPAR / postos locais (2024) \textbf{Tipo de
localidade:} Interior

{\def\LTcaptype{none} % do not increment counter
\begin{longtable}[]{@{}lll@{}}
\toprule\noalign{}
Combustível & USD/litro & Gs/litro (aprox.) \\
\midrule\noalign{}
\endhead
\bottomrule\noalign{}
\endlastfoot
Gasolina 93 oct & 0.98 & 7,252 \\
Gasolina 97 oct (premium) & 1.08 & 7,992 \\
Diesel & 0.91 & 6,734 \\
\end{longtable}
}

\begin{quote}
Preços podem variar ±5\% conforme posto e sazonalidade. Chaco e interior
remoto apresentam maior variação.
\end{quote}

\subsubsection{Cobertura Celular}

\textbf{Fonte:} CONATEL PY / operadoras (2024)

{\def\LTcaptype{none} % do not increment counter
\begin{longtable}[]{@{}ll@{}}
\toprule\noalign{}
Parâmetro & Valor \\
\midrule\noalign{}
\endhead
\bottomrule\noalign{}
\endlastfoot
Cobertura 4G população (dept.) & 92\% \\
Cobertura 4G área rural & 82\% \\
Melhor operadora & Tigo \\
Qualidade rural & boa \\
\end{longtable}
}

\begin{quote}
Para áreas rurais fora do núcleo urbano, recomenda-se chip Tigo como
principal e Personal como backup.
\end{quote}

\subsubsection{Internet}

\textbf{Fonte:} CONATEL / Speedtest Ookla (2024)

{\def\LTcaptype{none} % do not increment counter
\begin{longtable}[]{@{}ll@{}}
\toprule\noalign{}
Parâmetro & Valor \\
\midrule\noalign{}
\endhead
\bottomrule\noalign{}
\endlastfoot
Velocidade média download & 75 Mbps \\
Domicílios com internet (dept.) & 78\% \\
Tecnologia predominante & fibra/rádio \\
Opção rural & Starlink disponível (\textasciitilde USD 44/mês) \\
\end{longtable}
}

\subsubsection{Mercado Imobiliário e Terra
Rural}

\textbf{Fonte:} INDERT / Clasificados.com.py (2024)

{\def\LTcaptype{none} % do not increment counter
\begin{longtable}[]{@{}ll@{}}
\toprule\noalign{}
Tipo & Referência \\
\midrule\noalign{}
\endhead
\bottomrule\noalign{}
\endlastfoot
Terra agrícola alta prod. (USD/ha) & 6,900 \\
Imóvel urbano (USD/m²) & 1,300 \\
Aluguel 2 quartos (USD/mês) & 550 \\
\end{longtable}
}

\begin{quote}
Valores de referência departamental. Localidades menores podem ter
preços 20--40\% abaixo da capital departamental.
\end{quote}

\subsubsection{Saúde}

\textbf{Fonte:} MSPBS / IPS Paraguay (2024-2026), consolidação
departamental e proxy local

{\def\LTcaptype{none} % do not increment counter
\begin{longtable}[]{@{}ll@{}}
\toprule\noalign{}
Serviço & Disponibilidade \\
\midrule\noalign{}
\endhead
\bottomrule\noalign{}
\endlastfoot
USF / Posto de Saúde & sim \\
Hospital Regional & não \\
IPS (seguro social) & sim \\
Farmácia & sim \\
Distância ao hospital de referência & \textasciitilde18 km
(Encarnacion) \\
\end{longtable}
}

\textbf{Principais estabelecimentos:} USF local; referencia hospitalar
em Encarnacion

\textbf{Observação para imigrantes:} Atendimento primario local ou em
raio curto; casos de maior complexidade seguem para Encarnacion.
Cobertura privada continua recomendada para especialidades e urgências
de maior complexidade.
\newpage
\secaoDiagnostico{Carlos Antonio Lopez}{\mapaDistrito{0709}}{Carlos Antonio Lopez e um polo agricola de alta produtividade com
infraestrutura de saude em melhoria, porem penalizado pela segurança
publica fragil na regiao de fronteira.

\subsubsection{Por que sim}

\begin{itemize}
\tightlist
\item
  Solo de altissima qualidade e abundancia de agua.
\item
  Conectividade asfaltada estratégica para escoamento.
\end{itemize}

\subsubsection{Por que não}

\begin{itemize}
\tightlist
\item
  Historico recente de violencia armada e assaltos violentos.
\item
  Grande dispersao rural dificulta resposta policial rapida.
\end{itemize}

\subsubsection{Recomendação
Técnica}

Ideal para investimentos agricolas mecanizados, mas exige alto
investimento privado em segurança e protocolos de defesa civil. GSS 6.2.

\subsubsection{Analise de Sensibilidade}

\begin{itemize}
\tightlist
\item
  Otimista (Melhora Segurança E+2): GSS 6.6
\item
  Conservador (Crise Logistica A-1, B-1): GSS 5.8
\end{itemize}

\subsubsection{}

\begin{itemize}
\tightlist
\item
  \begin{enumerate}
  \def\labelenumi{\arabic{enumi})}
  \tightlist
  \item
    Censo INE 2022.
  \end{enumerate}
\item
  \begin{enumerate}
  \def\labelenumi{\arabic{enumi})}
  \setcounter{enumi}{1}
  \tightlist
  \item
    Relatorios de segurança Itapua Noticias (jan/2026).
  \end{enumerate}
\item
  \begin{enumerate}
  \def\labelenumi{\arabic{enumi})}
  \setcounter{enumi}{2}
  \tightlist
  \item
    Dados MOPC sobre Ruta PY07.
  \end{enumerate}
\item
  \begin{enumerate}
  \def\labelenumi{\arabic{enumi})}
  \setcounter{enumi}{3}
  \tightlist
  \item
    Inauguracao USF San Lorenzo (dez/2025).
  \end{enumerate}
\end{itemize}}{dist:07_Itapua:Carlos_Antonio_Lopez}
\begin{table}[ht!]
\small \begin{tabular}{p{3.0cm}p{0.8cm}p{6.0cm}} \toprule
\textbf{Critério} & \textbf{Nota} & \textbf{Justificativa} \\ \midrule
A - Ameaças Estratégicas & 6.0 & - \\
B - Risco Social & 6.5 & - \\
C - Riscos Naturais & 9.0 & - \\
D - Autossuficiência & 5.5 & - \\
E - Institucional & 5.0 & - \\
\midrule \textbf{GSS FINAL} & \textbf{6.2} & \textbf{Moderadamente Seguro (exige vigilancia em segurança).} \\ \bottomrule \end{tabular}
\end{table}
\subsubsection{Vulnerabilidades e
Mitigação}\label{vuln-Carlos_Antonio_Lopez-vulnerabilidades-e-mitigauxe7uxe3o}

\begin{itemize}
\item
  Segurança: Reforco de barreiras fisicas e sistemas de vigilancia em
  propriedades.
\item
  Logistica: Manter estoque de insumos para periodos de instabilidade na
  Rota PY07.
\item
  Saude: Plano de evacuacao para Encarnacion (Gran Hospital del Sur) em
  casos criticos.
\item
  Fonte INE (Censo 2022): 13.792 habitantes.
\item
  Fonte MOPC: Rota PY07 operacional e pavimentacao urbana recente.
\item
  Fonte MSPBS: Centro de Saude ampliado e nova USF em San Lorenzo
  (2025).
\end{itemize}
\subsubsection{Dossiê de Campo}\label{dossie-Carlos_Antonio_Lopez-dossiuxea-de-campo}
\begin{description}[style=nextline,leftmargin=0.5cm,font=\bfseries]
\item[Coordenadas]  26°24′S 54°45′W (geocodificado via Nominatim).
\end{description}
\textbf{Fonte climática:} NASA POWER Climatology API (período 2001-2020)

\textbf{Fonte luminosa:} estimativa world atlas

\paragraph{Irradiação Solar
(kWh/m²/dia)}\label{irradiauxe7uxe3o-solar-kwhmuxb2dia}

{\def\LTcaptype{none} % do not increment counter
\begin{longtable}[]{@{}
  >{\raggedright\arraybackslash}p{(\linewidth - 24\tabcolsep) * \real{0.0746}}
  >{\raggedright\arraybackslash}p{(\linewidth - 24\tabcolsep) * \real{0.0746}}
  >{\raggedright\arraybackslash}p{(\linewidth - 24\tabcolsep) * \real{0.0746}}
  >{\raggedright\arraybackslash}p{(\linewidth - 24\tabcolsep) * \real{0.0746}}
  >{\raggedright\arraybackslash}p{(\linewidth - 24\tabcolsep) * \real{0.0746}}
  >{\raggedright\arraybackslash}p{(\linewidth - 24\tabcolsep) * \real{0.0746}}
  >{\raggedright\arraybackslash}p{(\linewidth - 24\tabcolsep) * \real{0.0746}}
  >{\raggedright\arraybackslash}p{(\linewidth - 24\tabcolsep) * \real{0.0746}}
  >{\raggedright\arraybackslash}p{(\linewidth - 24\tabcolsep) * \real{0.0746}}
  >{\raggedright\arraybackslash}p{(\linewidth - 24\tabcolsep) * \real{0.0746}}
  >{\raggedright\arraybackslash}p{(\linewidth - 24\tabcolsep) * \real{0.0746}}
  >{\raggedright\arraybackslash}p{(\linewidth - 24\tabcolsep) * \real{0.0746}}
  >{\raggedright\arraybackslash}p{(\linewidth - 24\tabcolsep) * \real{0.1045}}@{}}
\toprule\noalign{}
\begin{minipage}[b]{\linewidth}\raggedright
Jan
\end{minipage} & \begin{minipage}[b]{\linewidth}\raggedright
Fev
\end{minipage} & \begin{minipage}[b]{\linewidth}\raggedright
Mar
\end{minipage} & \begin{minipage}[b]{\linewidth}\raggedright
Abr
\end{minipage} & \begin{minipage}[b]{\linewidth}\raggedright
Mai
\end{minipage} & \begin{minipage}[b]{\linewidth}\raggedright
Jun
\end{minipage} & \begin{minipage}[b]{\linewidth}\raggedright
Jul
\end{minipage} & \begin{minipage}[b]{\linewidth}\raggedright
Ago
\end{minipage} & \begin{minipage}[b]{\linewidth}\raggedright
Set
\end{minipage} & \begin{minipage}[b]{\linewidth}\raggedright
Out
\end{minipage} & \begin{minipage}[b]{\linewidth}\raggedright
Nov
\end{minipage} & \begin{minipage}[b]{\linewidth}\raggedright
Dez
\end{minipage} & \begin{minipage}[b]{\linewidth}\raggedright
Média
\end{minipage} \\
\midrule\noalign{}
\endhead
\bottomrule\noalign{}
\endlastfoot
6.47 & 6.02 & 5.33 & 4.39 & 3.24 & 2.80 & 3.13 & 3.90 & 4.47 & 5.26 &
6.32 & 6.54 & \textbf{4.82} \\
\end{longtable}
}

\textbf{Inclinação solar recomendada:} 26° N (anual) · 36° N (inverno
jun-ago) · 16° N (verão nov-jan)

\paragraph{Precipitação (mm/mês)}\label{precipitauxe7uxe3o-mmmuxeas}

{\def\LTcaptype{none} % do not increment counter
\begin{longtable}[]{@{}
  >{\raggedright\arraybackslash}p{(\linewidth - 24\tabcolsep) * \real{0.0704}}
  >{\raggedright\arraybackslash}p{(\linewidth - 24\tabcolsep) * \real{0.0704}}
  >{\raggedright\arraybackslash}p{(\linewidth - 24\tabcolsep) * \real{0.0704}}
  >{\raggedright\arraybackslash}p{(\linewidth - 24\tabcolsep) * \real{0.0704}}
  >{\raggedright\arraybackslash}p{(\linewidth - 24\tabcolsep) * \real{0.0704}}
  >{\raggedright\arraybackslash}p{(\linewidth - 24\tabcolsep) * \real{0.0704}}
  >{\raggedright\arraybackslash}p{(\linewidth - 24\tabcolsep) * \real{0.0704}}
  >{\raggedright\arraybackslash}p{(\linewidth - 24\tabcolsep) * \real{0.0704}}
  >{\raggedright\arraybackslash}p{(\linewidth - 24\tabcolsep) * \real{0.0704}}
  >{\raggedright\arraybackslash}p{(\linewidth - 24\tabcolsep) * \real{0.0704}}
  >{\raggedright\arraybackslash}p{(\linewidth - 24\tabcolsep) * \real{0.0704}}
  >{\raggedright\arraybackslash}p{(\linewidth - 24\tabcolsep) * \real{0.0704}}
  >{\raggedright\arraybackslash}p{(\linewidth - 24\tabcolsep) * \real{0.1549}}@{}}
\toprule\noalign{}
\begin{minipage}[b]{\linewidth}\raggedright
Jan
\end{minipage} & \begin{minipage}[b]{\linewidth}\raggedright
Fev
\end{minipage} & \begin{minipage}[b]{\linewidth}\raggedright
Mar
\end{minipage} & \begin{minipage}[b]{\linewidth}\raggedright
Abr
\end{minipage} & \begin{minipage}[b]{\linewidth}\raggedright
Mai
\end{minipage} & \begin{minipage}[b]{\linewidth}\raggedright
Jun
\end{minipage} & \begin{minipage}[b]{\linewidth}\raggedright
Jul
\end{minipage} & \begin{minipage}[b]{\linewidth}\raggedright
Ago
\end{minipage} & \begin{minipage}[b]{\linewidth}\raggedright
Set
\end{minipage} & \begin{minipage}[b]{\linewidth}\raggedright
Out
\end{minipage} & \begin{minipage}[b]{\linewidth}\raggedright
Nov
\end{minipage} & \begin{minipage}[b]{\linewidth}\raggedright
Dez
\end{minipage} & \begin{minipage}[b]{\linewidth}\raggedright
Total/ano
\end{minipage} \\
\midrule\noalign{}
\endhead
\bottomrule\noalign{}
\endlastfoot
149 & 129 & 134 & 174 & 156 & 118 & 94 & 76 & 113 & 206 & 180 & 182 &
\textbf{1712 mm} \\
\end{longtable}
}

\paragraph{Poluição Luminosa}\label{poluiuxe7uxe3o-luminosa}

{\def\LTcaptype{none} % do not increment counter
\begin{longtable}[]{@{}ll@{}}
\toprule\noalign{}
Parâmetro & Valor \\
\midrule\noalign{}
\endhead
\bottomrule\noalign{}
\endlastfoot
Escala Bortle & 5 --- Céu suburbano \\
Radiância artificial & 5.0 nW/cm²/sr \\
\end{longtable}
}
\paragraph{Solo (SoilGrids 2.0, média ponderada 0--30
cm)}

{\def\LTcaptype{none} % do not increment counter
\begin{longtable}[]{@{}ll@{}}
\toprule\noalign{}
Parâmetro & Valor \\
\midrule\noalign{}
\endhead
\bottomrule\noalign{}
\endlastfoot
pH (H₂O) & 5.2 \\
Carbono orgânico (SOC) & 18.23 g/kg \\
Argila & 47.57 \% \\
Areia & 24.08 \% \\
Silte (calc.) & 28.4 \% \\
Densidade aparente & 1.23 g/cm³ \\
\textbf{Aptidão agrícola} & \textbf{Média} \\
\end{longtable}
}

Fonte: ISRIC SoilGrids 2.0 via WCS. Coords: -26.4°, -54.75°. Média
ponderada camadas 0-5, 5-15, 15-30 cm.
\subsubsection{Indicadores Sociais}

\textbf{Fontes:} PNUD 2020, INE Censo 2022, INE EPHC 2023, Ministerio
Público 2024\\
\textbf{Nota:} Valores marcados como \emph{(dept.)} referem-se ao
departamento de Itapúa; valores \emph{(dist.)} são específicos deste
distrito. Dados marcados \emph{(est.)} são estimativas.

{\def\LTcaptype{none} % do not increment counter
\begin{longtable}[]{@{}
  >{\raggedright\arraybackslash}p{(\linewidth - 4\tabcolsep) * \real{0.4231}}
  >{\raggedright\arraybackslash}p{(\linewidth - 4\tabcolsep) * \real{0.2692}}
  >{\raggedright\arraybackslash}p{(\linewidth - 4\tabcolsep) * \real{0.3077}}@{}}
\toprule\noalign{}
\begin{minipage}[b]{\linewidth}\raggedright
Indicador
\end{minipage} & \begin{minipage}[b]{\linewidth}\raggedright
Valor
\end{minipage} & \begin{minipage}[b]{\linewidth}\raggedright
Âmbito
\end{minipage} \\
\midrule\noalign{}
\endhead
\bottomrule\noalign{}
\endlastfoot
IDH (2020) & 0,728 & dept. \\
Ranking IDH nacional & 4/18 & dept. \\
Esperança de vida & 74,0 anos & dept. \\
Escolaridade média & 7.4 anos & dist. \\
RNB per capita (USD PPA) & 10.100 & dept. \\
Pobreza monetária (\%) & 22,1\% & dept. \\
Pobreza extrema (\%) & 4,8\% & dept. \\
Índice de Gini & 0,432 & dept. \\
Acesso a água potável (\%) & 84,6\% & dept. \\
Acesso a saneamento (\%) & 83,1\% & dept. \\
Taxa de homicídios (est., /100k hab) & \textasciitilde4--6 (estimativa,
próxima à média nacional) & dept. \\
Índice de segurança & médio & dept. \\
População (Censo 2022) & 13.792 hab. & dist. \\
Idade mediana & 27.0 anos & dist. \\
Taxa de fecundidade & N/D & dist. \\
\end{longtable}
}
\subsubsection{Combustível}

\textbf{Referência:} PETROPAR / postos locais (2024) \textbf{Tipo de
localidade:} Interior

{\def\LTcaptype{none} % do not increment counter
\begin{longtable}[]{@{}lll@{}}
\toprule\noalign{}
Combustível & USD/litro & Gs/litro (aprox.) \\
\midrule\noalign{}
\endhead
\bottomrule\noalign{}
\endlastfoot
Gasolina 93 oct & 0.98 & 7,252 \\
Gasolina 97 oct (premium) & 1.08 & 7,992 \\
Diesel & 0.91 & 6,734 \\
\end{longtable}
}

\begin{quote}
Preços podem variar ±5\% conforme posto e sazonalidade. Chaco e interior
remoto apresentam maior variação.
\end{quote}

\subsubsection{Cobertura Celular}

\textbf{Fonte:} CONATEL PY / operadoras (2024)

{\def\LTcaptype{none} % do not increment counter
\begin{longtable}[]{@{}ll@{}}
\toprule\noalign{}
Parâmetro & Valor \\
\midrule\noalign{}
\endhead
\bottomrule\noalign{}
\endlastfoot
Cobertura 4G população (dept.) & 92\% \\
Cobertura 4G área rural & 82\% \\
Melhor operadora & Tigo \\
Qualidade rural & boa \\
\end{longtable}
}

\begin{quote}
Para áreas rurais fora do núcleo urbano, recomenda-se chip Tigo como
principal e Personal como backup.
\end{quote}

\subsubsection{Internet}

\textbf{Fonte:} CONATEL / Speedtest Ookla (2024)

{\def\LTcaptype{none} % do not increment counter
\begin{longtable}[]{@{}ll@{}}
\toprule\noalign{}
Parâmetro & Valor \\
\midrule\noalign{}
\endhead
\bottomrule\noalign{}
\endlastfoot
Velocidade média download & 75 Mbps \\
Domicílios com internet (dept.) & 78\% \\
Tecnologia predominante & fibra/rádio \\
Opção rural & Starlink disponível (\textasciitilde USD 44/mês) \\
\end{longtable}
}

\subsubsection{Mercado Imobiliário e Terra
Rural}

\textbf{Fonte:} INDERT / Clasificados.com.py (2024)

{\def\LTcaptype{none} % do not increment counter
\begin{longtable}[]{@{}ll@{}}
\toprule\noalign{}
Tipo & Referência \\
\midrule\noalign{}
\endhead
\bottomrule\noalign{}
\endlastfoot
Terra agrícola alta prod. (USD/ha) & 6,900 \\
Imóvel urbano (USD/m²) & 1,300 \\
Aluguel 2 quartos (USD/mês) & 550 \\
\end{longtable}
}

\begin{quote}
Valores de referência departamental. Localidades menores podem ter
preços 20--40\% abaixo da capital departamental.
\end{quote}

\subsubsection{Saúde}

\textbf{Fonte:} MSPBS / IPS Paraguay (2024-2026), consolidação
departamental e proxy local

{\def\LTcaptype{none} % do not increment counter
\begin{longtable}[]{@{}ll@{}}
\toprule\noalign{}
Serviço & Disponibilidade \\
\midrule\noalign{}
\endhead
\bottomrule\noalign{}
\endlastfoot
USF / Posto de Saúde & sim \\
Hospital Regional & sim \\
IPS (seguro social) & sim \\
Farmácia & sim \\
Distância ao hospital de referência & \textasciitilde190 km
(Encarnacion) \\
\end{longtable}
}

\textbf{Principais estabelecimentos:} USF local; referencia hospitalar
em Encarnacion

\textbf{Observação para imigrantes:} Atendimento primario local ou em
raio curto; casos de maior complexidade seguem para Encarnacion.
Cobertura privada continua recomendada para especialidades e urgências
de maior complexidade.
\newpage
\secaoDiagnostico{Carmen del Parana}{\mapaDistrito{0707}}{Carmen del Parana combina alta qualidade de vida, infraestrutura moderna
e seguranca acima da media regional, sendo um dos distritos mais
promissores de Itapua.

\subsubsection{Por que sim}

\begin{itemize}
\tightlist
\item
  Infraestrutura de transporte e lazer de primeiro nivel.
\item
  Proximidade estrategica com o maior centro regional (Encarnacion).
\item
  Solo altamente produtivo para graos.
\end{itemize}

\subsubsection{Por que não}

\begin{itemize}
\tightlist
\item
  Vulnerabilidade a crimes de fronteira fluvial.
\item
  Populacao reduzida limita a escala de defesa civil autonoma.
\end{itemize}

\subsubsection{Recomendação
Técnica}

Altamente recomendado para realocacao residencial e investimentos em
agro-turismo. GSS 6.6.

\subsubsection{Analise de Sensibilidade}

\begin{itemize}
\tightlist
\item
  Otimista (Reducao Crimes E+2): GSS 7.0
\item
  Conservador (Degradacao da PY01 A-2): GSS 6.1
\end{itemize}

\subsubsection{}

\begin{itemize}
\tightlist
\item
  \begin{enumerate}
  \def\labelenumi{\arabic{enumi})}
  \tightlist
  \item
    Censo INE 2022.
  \end{enumerate}
\item
  \begin{enumerate}
  \def\labelenumi{\arabic{enumi})}
  \setcounter{enumi}{1}
  \tightlist
  \item
    Relatorios de seguranca Itapua Noticias (2025).
  \end{enumerate}
\item
  \begin{enumerate}
  \def\labelenumi{\arabic{enumi})}
  \setcounter{enumi}{2}
  \tightlist
  \item
    Dados MOPC sobre Corredor Graneros del Sur.
  \end{enumerate}
\item
  \begin{enumerate}
  \def\labelenumi{\arabic{enumi})}
  \setcounter{enumi}{3}
  \tightlist
  \item
    Fortalecimento do Centro de Saude (MSPBS 2025).
  \end{enumerate}
\end{itemize}}{dist:07_Itapua:Carmen_del_Parana}
\begin{table}[ht!]
\small \begin{tabular}{p{3.0cm}p{0.8cm}p{6.0cm}} \toprule
\textbf{Critério} & \textbf{Nota} & \textbf{Justificativa} \\ \midrule
A - Ameaças Estratégicas & 7.0 & - \\
B - Risco Social & 7.0 & - \\
C - Riscos Naturais & 9.0 & - \\
D - Autossuficiência & 6.0 & - \\
E - Institucional & 4.5 & - \\
\midrule \textbf{GSS FINAL} & \textbf{6.6} & \textbf{Seguro (alto potencial turistico e agricola).} \\ \bottomrule \end{tabular}
\end{table}
\subsubsection{Vulnerabilidades e
Mitigação}\label{vuln-Carmen_del_Parana-vulnerabilidades-e-mitigauxe7uxe3o}

\begin{itemize}
\item
  Segurança Fluvial: Monitoramento de acessos via Rio Parana para
  prevenir incursões criminosas.
\item
  Dependencia Turistica: Diversificar a economia local para alem do
  veraneio.
\item
  Saude: Embora o Centro de Saude local tenha sido fortalecido (2025),
  casos graves exigem deslocamento para Encarnacion.
\item
  Fonte INE (Censo 2022): 6.417 habitantes.
\item
  Fonte MOPC: Corredor Graneros del Sur e Ruta PY01 em excelente estado.
\item
  Fonte MSPBS: Centro de Saude com novos equipamentos e medicos
  residentes (2025).
\end{itemize}
\subsubsection{Dossiê de Campo}\label{dossie-Carmen_del_Parana-dossiuxea-de-campo}
\begin{description}[style=nextline,leftmargin=0.5cm,font=\bfseries]
\item[Coordenadas]  27°10′S 56°04′W (geocodificado via Nominatim).
\end{description}
\textbf{Fonte climática:} NASA POWER Climatology API (período 2001-2020)

\textbf{Fonte luminosa:} estimativa world atlas

\paragraph{Irradiação Solar
(kWh/m²/dia)}\label{irradiauxe7uxe3o-solar-kwhmuxb2dia}

{\def\LTcaptype{none} % do not increment counter
\begin{longtable}[]{@{}
  >{\raggedright\arraybackslash}p{(\linewidth - 24\tabcolsep) * \real{0.0746}}
  >{\raggedright\arraybackslash}p{(\linewidth - 24\tabcolsep) * \real{0.0746}}
  >{\raggedright\arraybackslash}p{(\linewidth - 24\tabcolsep) * \real{0.0746}}
  >{\raggedright\arraybackslash}p{(\linewidth - 24\tabcolsep) * \real{0.0746}}
  >{\raggedright\arraybackslash}p{(\linewidth - 24\tabcolsep) * \real{0.0746}}
  >{\raggedright\arraybackslash}p{(\linewidth - 24\tabcolsep) * \real{0.0746}}
  >{\raggedright\arraybackslash}p{(\linewidth - 24\tabcolsep) * \real{0.0746}}
  >{\raggedright\arraybackslash}p{(\linewidth - 24\tabcolsep) * \real{0.0746}}
  >{\raggedright\arraybackslash}p{(\linewidth - 24\tabcolsep) * \real{0.0746}}
  >{\raggedright\arraybackslash}p{(\linewidth - 24\tabcolsep) * \real{0.0746}}
  >{\raggedright\arraybackslash}p{(\linewidth - 24\tabcolsep) * \real{0.0746}}
  >{\raggedright\arraybackslash}p{(\linewidth - 24\tabcolsep) * \real{0.0746}}
  >{\raggedright\arraybackslash}p{(\linewidth - 24\tabcolsep) * \real{0.1045}}@{}}
\toprule\noalign{}
\begin{minipage}[b]{\linewidth}\raggedright
Jan
\end{minipage} & \begin{minipage}[b]{\linewidth}\raggedright
Fev
\end{minipage} & \begin{minipage}[b]{\linewidth}\raggedright
Mar
\end{minipage} & \begin{minipage}[b]{\linewidth}\raggedright
Abr
\end{minipage} & \begin{minipage}[b]{\linewidth}\raggedright
Mai
\end{minipage} & \begin{minipage}[b]{\linewidth}\raggedright
Jun
\end{minipage} & \begin{minipage}[b]{\linewidth}\raggedright
Jul
\end{minipage} & \begin{minipage}[b]{\linewidth}\raggedright
Ago
\end{minipage} & \begin{minipage}[b]{\linewidth}\raggedright
Set
\end{minipage} & \begin{minipage}[b]{\linewidth}\raggedright
Out
\end{minipage} & \begin{minipage}[b]{\linewidth}\raggedright
Nov
\end{minipage} & \begin{minipage}[b]{\linewidth}\raggedright
Dez
\end{minipage} & \begin{minipage}[b]{\linewidth}\raggedright
Média
\end{minipage} \\
\midrule\noalign{}
\endhead
\bottomrule\noalign{}
\endlastfoot
6.74 & 6.25 & 5.37 & 4.27 & 3.19 & 2.70 & 3.05 & 3.74 & 4.44 & 5.36 &
6.49 & 6.87 & \textbf{4.87} \\
\end{longtable}
}

\textbf{Inclinação solar recomendada:} 27° N (anual) · 37° N (inverno
jun-ago) · 17° N (verão nov-jan)

\paragraph{Precipitação (mm/mês)}\label{precipitauxe7uxe3o-mmmuxeas}

{\def\LTcaptype{none} % do not increment counter
\begin{longtable}[]{@{}
  >{\raggedright\arraybackslash}p{(\linewidth - 24\tabcolsep) * \real{0.0704}}
  >{\raggedright\arraybackslash}p{(\linewidth - 24\tabcolsep) * \real{0.0704}}
  >{\raggedright\arraybackslash}p{(\linewidth - 24\tabcolsep) * \real{0.0704}}
  >{\raggedright\arraybackslash}p{(\linewidth - 24\tabcolsep) * \real{0.0704}}
  >{\raggedright\arraybackslash}p{(\linewidth - 24\tabcolsep) * \real{0.0704}}
  >{\raggedright\arraybackslash}p{(\linewidth - 24\tabcolsep) * \real{0.0704}}
  >{\raggedright\arraybackslash}p{(\linewidth - 24\tabcolsep) * \real{0.0704}}
  >{\raggedright\arraybackslash}p{(\linewidth - 24\tabcolsep) * \real{0.0704}}
  >{\raggedright\arraybackslash}p{(\linewidth - 24\tabcolsep) * \real{0.0704}}
  >{\raggedright\arraybackslash}p{(\linewidth - 24\tabcolsep) * \real{0.0704}}
  >{\raggedright\arraybackslash}p{(\linewidth - 24\tabcolsep) * \real{0.0704}}
  >{\raggedright\arraybackslash}p{(\linewidth - 24\tabcolsep) * \real{0.0704}}
  >{\raggedright\arraybackslash}p{(\linewidth - 24\tabcolsep) * \real{0.1549}}@{}}
\toprule\noalign{}
\begin{minipage}[b]{\linewidth}\raggedright
Jan
\end{minipage} & \begin{minipage}[b]{\linewidth}\raggedright
Fev
\end{minipage} & \begin{minipage}[b]{\linewidth}\raggedright
Mar
\end{minipage} & \begin{minipage}[b]{\linewidth}\raggedright
Abr
\end{minipage} & \begin{minipage}[b]{\linewidth}\raggedright
Mai
\end{minipage} & \begin{minipage}[b]{\linewidth}\raggedright
Jun
\end{minipage} & \begin{minipage}[b]{\linewidth}\raggedright
Jul
\end{minipage} & \begin{minipage}[b]{\linewidth}\raggedright
Ago
\end{minipage} & \begin{minipage}[b]{\linewidth}\raggedright
Set
\end{minipage} & \begin{minipage}[b]{\linewidth}\raggedright
Out
\end{minipage} & \begin{minipage}[b]{\linewidth}\raggedright
Nov
\end{minipage} & \begin{minipage}[b]{\linewidth}\raggedright
Dez
\end{minipage} & \begin{minipage}[b]{\linewidth}\raggedright
Total/ano
\end{minipage} \\
\midrule\noalign{}
\endhead
\bottomrule\noalign{}
\endlastfoot
150 & 123 & 146 & 184 & 148 & 107 & 85 & 80 & 111 & 206 & 194 & 179 &
\textbf{1713 mm} \\
\end{longtable}
}

\paragraph{Poluição Luminosa}\label{poluiuxe7uxe3o-luminosa}

{\def\LTcaptype{none} % do not increment counter
\begin{longtable}[]{@{}ll@{}}
\toprule\noalign{}
Parâmetro & Valor \\
\midrule\noalign{}
\endhead
\bottomrule\noalign{}
\endlastfoot
Escala Bortle & 5 --- Céu suburbano \\
Radiância artificial & 5.0 nW/cm²/sr \\
\end{longtable}
}
\paragraph{Solo (SoilGrids 2.0, média ponderada 0--30
cm)}

{\def\LTcaptype{none} % do not increment counter
\begin{longtable}[]{@{}ll@{}}
\toprule\noalign{}
Parâmetro & Valor \\
\midrule\noalign{}
\endhead
\bottomrule\noalign{}
\endlastfoot
pH (H₂O) & 5.2 \\
Carbono orgânico (SOC) & 20.0 g/kg \\
Argila & 42.05 \% \\
Areia & 22.47 \% \\
Silte (calc.) & 35.5 \% \\
Densidade aparente & 1.24 g/cm³ \\
\textbf{Aptidão agrícola} & \textbf{Média} \\
\end{longtable}
}

Fonte: ISRIC SoilGrids 2.0 via WCS. Coords: -27.1667°, -56.0667°. Média
ponderada camadas 0-5, 5-15, 15-30 cm.
\subsubsection{Indicadores Sociais}

\textbf{Fontes:} PNUD 2020, INE Censo 2022, INE EPHC 2023, Ministerio
Público 2024\\
\textbf{Nota:} Valores marcados como \emph{(dept.)} referem-se ao
departamento de Itapúa; valores \emph{(dist.)} são específicos deste
distrito. Dados marcados \emph{(est.)} são estimativas.

{\def\LTcaptype{none} % do not increment counter
\begin{longtable}[]{@{}
  >{\raggedright\arraybackslash}p{(\linewidth - 4\tabcolsep) * \real{0.4231}}
  >{\raggedright\arraybackslash}p{(\linewidth - 4\tabcolsep) * \real{0.2692}}
  >{\raggedright\arraybackslash}p{(\linewidth - 4\tabcolsep) * \real{0.3077}}@{}}
\toprule\noalign{}
\begin{minipage}[b]{\linewidth}\raggedright
Indicador
\end{minipage} & \begin{minipage}[b]{\linewidth}\raggedright
Valor
\end{minipage} & \begin{minipage}[b]{\linewidth}\raggedright
Âmbito
\end{minipage} \\
\midrule\noalign{}
\endhead
\bottomrule\noalign{}
\endlastfoot
IDH (2020) & 0,728 & dept. \\
Ranking IDH nacional & 4/18 & dept. \\
Esperança de vida & 74,0 anos & dept. \\
Escolaridade média & 8.5 anos & dist. \\
RNB per capita (USD PPA) & 10.100 & dept. \\
Pobreza monetária (\%) & 22,1\% & dept. \\
Pobreza extrema (\%) & 4,8\% & dept. \\
Índice de Gini & 0,432 & dept. \\
Acesso a água potável (\%) & 84,6\% & dept. \\
Acesso a saneamento (\%) & 83,1\% & dept. \\
Taxa de homicídios (est., /100k hab) & \textasciitilde4--6 (estimativa,
próxima à média nacional) & dept. \\
Índice de segurança & médio & dept. \\
População (Censo 2022) & 6.417 hab. & dist. \\
Idade mediana & 31.0 anos & dist. \\
Taxa de fecundidade & N/D & dist. \\
\end{longtable}
}
\subsubsection{Combustível}

\textbf{Referência:} PETROPAR / postos locais (2024) \textbf{Tipo de
localidade:} Interior

{\def\LTcaptype{none} % do not increment counter
\begin{longtable}[]{@{}lll@{}}
\toprule\noalign{}
Combustível & USD/litro & Gs/litro (aprox.) \\
\midrule\noalign{}
\endhead
\bottomrule\noalign{}
\endlastfoot
Gasolina 93 oct & 0.98 & 7,252 \\
Gasolina 97 oct (premium) & 1.08 & 7,992 \\
Diesel & 0.91 & 6,734 \\
\end{longtable}
}

\begin{quote}
Preços podem variar ±5\% conforme posto e sazonalidade. Chaco e interior
remoto apresentam maior variação.
\end{quote}

\subsubsection{Cobertura Celular}

\textbf{Fonte:} CONATEL PY / operadoras (2024)

{\def\LTcaptype{none} % do not increment counter
\begin{longtable}[]{@{}ll@{}}
\toprule\noalign{}
Parâmetro & Valor \\
\midrule\noalign{}
\endhead
\bottomrule\noalign{}
\endlastfoot
Cobertura 4G população (dept.) & 92\% \\
Cobertura 4G área rural & 82\% \\
Melhor operadora & Tigo \\
Qualidade rural & boa \\
\end{longtable}
}

\begin{quote}
Para áreas rurais fora do núcleo urbano, recomenda-se chip Tigo como
principal e Personal como backup.
\end{quote}

\subsubsection{Internet}

\textbf{Fonte:} CONATEL / Speedtest Ookla (2024)

{\def\LTcaptype{none} % do not increment counter
\begin{longtable}[]{@{}ll@{}}
\toprule\noalign{}
Parâmetro & Valor \\
\midrule\noalign{}
\endhead
\bottomrule\noalign{}
\endlastfoot
Velocidade média download & 75 Mbps \\
Domicílios com internet (dept.) & 78\% \\
Tecnologia predominante & fibra/rádio \\
Opção rural & Starlink disponível (\textasciitilde USD 44/mês) \\
\end{longtable}
}

\subsubsection{Mercado Imobiliário e Terra
Rural}

\textbf{Fonte:} INDERT / Clasificados.com.py (2024)

{\def\LTcaptype{none} % do not increment counter
\begin{longtable}[]{@{}ll@{}}
\toprule\noalign{}
Tipo & Referência \\
\midrule\noalign{}
\endhead
\bottomrule\noalign{}
\endlastfoot
Terra agrícola alta prod. (USD/ha) & 6,900 \\
Imóvel urbano (USD/m²) & 1,300 \\
Aluguel 2 quartos (USD/mês) & 550 \\
\end{longtable}
}

\begin{quote}
Valores de referência departamental. Localidades menores podem ter
preços 20--40\% abaixo da capital departamental.
\end{quote}

\subsubsection{Saúde}

\textbf{Fonte:} MSPBS / IPS Paraguay (2024-2026), consolidação
departamental e proxy local

{\def\LTcaptype{none} % do not increment counter
\begin{longtable}[]{@{}ll@{}}
\toprule\noalign{}
Serviço & Disponibilidade \\
\midrule\noalign{}
\endhead
\bottomrule\noalign{}
\endlastfoot
USF / Posto de Saúde & sim \\
Hospital Regional & sim \\
IPS (seguro social) & sim \\
Farmácia & sim \\
Distância ao hospital de referência & \textasciitilde38 km
(Encarnacion) \\
\end{longtable}
}

\textbf{Principais estabelecimentos:} USF local; referencia hospitalar
em Encarnacion

\textbf{Observação para imigrantes:} Atendimento primario local ou em
raio curto; casos de maior complexidade seguem para Encarnacion.
Cobertura privada continua recomendada para especialidades e urgências
de maior complexidade.
\newpage
